\chapter{Loi générale des masses : échelle absolue d’action \texorpdfstring{\(10^{-34}\)}{10⁻³⁴}, domaine \texorpdfstring{\(81/56/13/324\)}{81/56/13/324}, secteur nul et construction opératorielle du spectre matériel}
\label{chap:ley-general-masas-dominio-completo}

\section{Domaine opératoriel de la loi générale des masses : 81 cellules, 56 familles, 13 multisections, 324 parcours et secteur nul}
La loi générale des masses agit sur les sections persistantes de l’atlas orienté, et non sur un tableau d’espèces. Son domaine interne contient \(81\) cellules, \(56\) familles de trajectoires, \(13\) multisections et \(324\) parcours orientés, outre une carte nulle séparée. Les recensements de sorties physiques sont ultérieurs et ne modifient pas ce domaine.

\section{Origine déterminantale de l’échelle absolue}
La forme entière \(H_4\) a pour forme normale de Smith \(\operatorname{diag}(1,2,2,4)\), si bien que son conoyau est d’ordre \(16\). Il en résulte
\[
 \Sigma_H=\varphi_{\rm HMT}\alpha_{\rm HMT}^{16},
 \qquad \operatorname{ord}_{10}(\Sigma_H)=-34.
\]
La chaîne absolue est
\[
 \Sigma_H\to\hbar_{\rm pre/ret}\to\omega_0
 \to\varepsilon_{\rm clk}\to m_{\rm clk}
 \to\widehat M^{(0)}\to\widehat M^{(\beta)}.
\]
La décennie \(-34\) n’est jamais introduite à partir du système SI ni absorbée dans le correcteur relatif.

\section{Construction directrice de la loi et du domaine matériel}
Le développement intégral de la loi est réparti ici par résultat : état, parcours, registre, signature, caractère, tours, branche absolue et opérateur. Son ordre interne est conservé, tandis que les réalisations par espèce sont réservées au chapitre correspondant.

\begingroup
\let\chapter\section
\let\section\subsection
\let\subsection\subsubsection
\let\subsubsection\paragraph
\let\clearpage\relax
\section{Définition spectrale de la masse comme mesure de persistance de \(1\) parcours couplé à un régime physique}
La masse est la mesure spectrale de la persistance d'un parcours lorsqu'une section généalogique se couple à un régime physique. Elle n'est pas un attribut scalaire attaché d'emblée à un nom de particule. La dynamique APP--TRIT--TPK construit le parcours, sa mémoire, sa signature et son caractère ; la Mécanique Dimensionnelle réalise cette généalogie dans une droite de masse et détermine, par le couplage, le sous-espace observable et son spectre.

Une particule élémentaire est représentée par une section persistante ; une famille non centrale, par une multisection stable ; une particule composée, par une composition d'enregistrements généalogiques avec frontière de liaison ; une résonance, par un pôle de l'opérateur couplé ; un secteur topologique, par son anneau de fusion et sa représentation de tressage avant toute carte de masse ; une section virtuelle, par son horizon terminal de prolongement. Ces réalisations sont distinctes et ne doivent pas être aplaties en une unique liste de nombres.

La grandeur externe n'intervient jamais dans la sélection de la cellule, du parcours, de la signature, du lecteur, du coefficient de tour ou du vecteur propre. La confrontation métrologique est effectuée après la construction et la fixation du résultat interne.

\Needspace{7\baselineskip}
\section{Espace complet des états de masse}
Soit \(\mathcal C_{81}\) l'atlas cellulaire. Chaque cellule admet quatre orientations two--way, de sorte que l'espace fini des parcours orientés est

\[
 \mathcal R_{324}
 =\{(c,h_0,v_0):c\in\mathcal C_{81},\ (h_0,v_0)\in\{\pm1\}^2\}.
\]
La projection cellulaire, la relation de famille et la stratification grossière produisent

\[
 |\mathcal C_{81}|=81,\qquad
 |\mathcal F_{56}|=56,\qquad
 |\mathcal M_{13}|=13,\qquad
 |\mathcal R_{324}|=324.
\]
L'électron occupe l'unique point fixe central. Les douze multisections non centrales sont des fibres finies stables par inversion cellulaire et n'admettent généralement pas de section ponctuelle équivariante. Leur objet naturel est donc un opérateur sur la fibre, et non un nombre choisi parmi ses cellules. Le secteur nul est incorporé au moyen de la section zéro de la droite de masse, et non comme valeur d'un caractère positif.

Le catalogue physique résultant distingue trois leptons chargés, trois réalisations neutriniques d'interaction, six saveurs de quarks avec leurs orbites de couleur, deux réalisations de jauge nulles, deux bosons de jauge massifs, un secteur scalaire, deux familles composées, trois codomaines topologiques et un système de sections virtuelles. Cette classification ultérieure ne limite pas le nombre d'états composés ou résonants : composition et excitation engendrent des familles supplémentaires au sein de leurs codomaines respectifs.

\Needspace{7\baselineskip}
\section{Trajectoire généalogique d'une réalisation de masse}
La loi ne commence ni par une étiquette de particule ni par un nombre. Son ordre est

\[
\boxed{
\begin{aligned}
\mathrm{APP}\to\mathrm{TRIT}\to\mathrm{TPK}
&\to\widetilde\Gamma^{\rm surv}
\to\Gamma^{\rm obs}
\to L_{108}
\to\sigma\in\mathbb Z^5\\
&\to\chi_0
\to\chi_{\rm full}
\to\widehat{\mathcal M}
\to\text{couplage}\\
&\to\text{spectre/pôle}
\to\mathcal L_M
\to\text{confrontation}.
\end{aligned}}
\]
Ici :

\begin{itemize}
\item APP dispose l'incompatibilité locale, les feuillets somme/produit et leur orientation ;
\item le TRIT conserve signe, seuil et ouverture bidirectionnelle ;
\item le TPK fait évoluer phase, mémoire, retenue, survie et retour ;
\item une particule est une section persistante d'une extension, et son parcours observable est
l'histoire qu'elle laisse dans l'enregistrement généalogique ;

\item la signature entière est additive ;
\item le caractère convertit la composition additive en un rapport positif ;
\item les tours raffinent le caractère au moyen d'un unique opérateur harmonique global ;
\item la droite de masse fournit le type dimensionnel ;
\item mesurer, c'est se coupler : le couplage détermine quel sous-espace et quel spectre sont
physiquement accessibles, sans consulter la grandeur qui sera ensuite mesurée.

\end{itemize}
La masse ne sélectionne pas le parcours. Le parcours persistant, l'enregistrement généalogique et le couplage sélectionnent l'objet spectral dont la masse est lue.

\Needspace{7\baselineskip}
\subsection{Composantes conservées de l'état enrichi du TPK dans la trajectoire de masse}
L'état qui alimente la loi des masses est

\[
 X_K=(w_6,w_{30},R_{36},g_K,B_K,p_K,c_K,\varepsilon_K,\omega_K),
\]
où sont conservés mot court, mot prolongé, région, phase nonadique, mémoire de bloc, préfixe, retenue, feuillet et enroulement. La chaîne

\[
 w_6\longrightarrow w_{30}\longrightarrow R_{36}
 \longrightarrow G_9
\]
ne projette pas directement une masse : elle construit l'identité persistante qui sera ensuite lue.

\Needspace{7\baselineskip}
\subsection{Filtration nonadique avec mémoire}
À partir de chaque état, on obtient une fibre de raffinement formée de \(729\) sous-cylindres,

\[
 \{C_{K,j}:0\le j<729\}.
\]
L'élagage \(G_9\) satisfait

\[
 G_9(C_{K,j})\in\{0,1\},\qquad
 \sum_{j=0}^{728}G_9(C_{K,j})=1.
\]
Si \(j_K\) est le survivant,

\[
 P_{K+1}=729P_K+j_K.
\]
Après neuf pas,

\[
 g(K+9)=g(K),
\]
mais changent, en général,

\[
 (B_{K+9},P_{K+9},\operatorname{pref}_{K+9},
 \operatorname{Witt}_{K+9})
 \ne
 (B_K,P_K,\operatorname{pref}_K,\operatorname{Witt}_K).
\]
Telle est la forme équationnelle de la persistance : la phase revient, et non l'état. La masse peut reconnaître une particule après le retour parce que l'enregistrement généalogique conserve la différence.

\Needspace{7\baselineskip}
\subsection{Transduction de la retenue nonadique vers l'état signé dodécaphasique}
Pour une coordonnée temporelle \(\theta\), nous définissons

\[
 g_0(\theta)=\theta\bmod9,\qquad
 \beta(\theta)=
 \left\lfloor\frac{\widetilde\theta}{9}\right\rfloor\bmod12.
\]
Alors

\[
 g_0(\theta+9)=g_0(\theta),\qquad
 \beta(\theta+9)=\beta(\theta)+1\pmod{12}.
\]
Ainsi, le retour nonadique conserve la phase et avance d'un cran de mémoire. Les 108 temps sont

\[
 108=9\cdot12,
\]
de sorte que chaque enregistrement généalogique contient un cycle complet des deux échelles. Les 36.864 compositions de retenue décrivent les combinaisons locales de cette mémoire bidirectionnelle.

\Needspace{7\baselineskip}
\subsection{Complétion et conservation de l'unité}
La frontière en base mille se décompose comme

\[
 999=729+270=729+243+27,
\]
tandis que la complétion est

\[
 1000=729+271=729+243+27+1,
\qquad
 271=270+1.
\]
Le \(+1\) est l'unique sommet entièrement frontalier et réalise la retenue qui restitue l'unité. La partition \(729/271\) distingue noyau actif et bord médiateur ; elle n'est pas une approximation du nombre d'or. La même conservation permet à un parcours de changer de bloc sans perdre son identité normalisée.

\Needspace{7\baselineskip}
\subsection{Théorème de l'atlas orienté complet}
Chaque cellule \(c\in\mathcal C_{81}\) possède quatre orientations

\[
 (h_0,v_0)\in\{\pm1\}^2.
\]
L’application

\[
 (c,h_0,v_0)\longmapsto\Gamma(c,h_0,v_0)
\]
est bijectif sur \(\mathcal R_{324}\). En conséquence,

\[
 |\mathcal R_{324}|=4|\mathcal C_{81}|=324.
\]
Chaque parcours possède 108 entrées d'enregistrement généalogique, de sorte que le dénombrement temporel total est

\[
 81\cdot108=8748.
\]
Le quadruplet

\[
 (\text{ligne germe},\text{colonne de graines},h_0,v_0)
\]
détermine de manière unique le parcours, son mot \(w_6\), son mot dodécaphasique, ses retours, sa mémoire, sa retenue, sa signature historique et ses lecteurs. Les 324 parcours peuvent donc non seulement être énumérés, mais aussi reconstruits et réunis exactement, champ par champ, sans utiliser de masse.

\Needspace{7\baselineskip}
\subsection{Quotients de famille et de multisection}
Le quotient des 81 cellules par histoire de parcours produit 56 familles. Leurs multiplicités se répartissent comme

\[
 41\times1,\quad10\times2,\quad2\times3,\quad1\times4,\quad2\times5,
\]
et

\[
 41+20+6+4+10=81.
\]
La stratification en treize multisections a pour rangs

\[
 1,\quad2,\quad4,\quad6,\quad8,\quad12,\quad16
\]
avec multiplicités

\[
 1,\quad2,\quad3,\quad2,\quad3,\quad1,\quad1,
\]
et pour somme des rangs

\[
 1+4+12+12+24+12+16=81.
\]
L'électron occupe le bloc de rang un. Les autres secteurs se situent dans des blocs de rang supérieur et exigent un projecteur de représentation, et non une cellule nominale.

\Needspace{7\baselineskip}
\subsection{Signature, lecteurs et mémoire par parcours}
Pour chaque \(\Gamma\in\mathcal R_{324}\), on calcule :

\[
 \sigma_\Gamma\in\mathbb Z^5,\qquad
 K_D(\Gamma),\qquad K_\Omega(\Gamma),
\]
avec les premiers retours de classe, d'orbite, de cellule et d'état cinématique ; la séparation des préfixes ; la dette de retenue ; l'ambiguïté ; les triades stabilisées ; le feuillet et l'enroulement. La famille \(K_D\) présente 14 profils ; \(K_\Omega\), neuf ; ensemble, ils forment 40 paires et 18 coïncidences. Cette diversité constitue l'information fine qui permet de passer d'une fibre grossière à un opérateur physique.

\Needspace{7\baselineskip}
\section{Du défaut APP à la signature entière}
\Needspace{7\baselineskip}
\subsection{Défaut primordial et déploiement}
Au centre APP,

\[
 B_c=9P_++5P_-,\qquad 9-5=4.
\]
Le saut local primordial est 4. Le défaut global \(54\) appartient à un déploiement ultérieur \(2\to6\to54\) ; il ne doit pas être rétroprojeté sur le centre comme s'il s'agissait de la même opération.

\Needspace{7\baselineskip}
\subsection{\(12\) événements et diviseur dodécaphasique}
La coordonnée orientée ne reçoit pas arbitrairement un facteur \(1/6\). La chaîne est

\[
12\ \text{événements}
\to s_C^{\rm tot}
\to 1+5+6
\to s_C^{(12)}=\frac{s_C^{\rm tot}}6
\to W_\pm=6n_A\pm s_C
\to\operatorname{SNF}(1,12)
\to\chi_{\rm mass}.
\]
Le \(1/6\) est le quotient orienté de la clôture de douze événements. Ce n'est ni le résidu (-1/12), ni la pseudo-inverse d'une valence, ni l'incidence hexade--octade.

\Needspace{7\baselineskip}
\subsection{Construction de la signature entière complète à partir de l'enregistrement du parcours}
L'enregistrement généalogique CT--108 produit

\[
 \sigma(\Gamma)
 =\bigl(n_A,s_C,k_{\Delta_4},b_{120},b_\zeta\bigr)\in\mathbb Z^5.
\]
Les coordonnées ont des provenances distinctes :

\begin{itemize}
\item \(n_A\) : durée/persistance dominante ;
\item \(s_C\) : orientation nette de la bicouche, divisée par six après la clôture
dodécaphasique ;

\item \(k_{\Delta_4}\) : canal de torsion minimale ;
\item \(b_{120}\) : coefficient d'événement de l'enregistrement ;
\item \(b_\zeta\) : autocorrélation de pas trois.
\end{itemize}
On définit

\[
 \zeta_{270}:=135\Delta_4^2,
\]
et reste distinct de

\[
 s_{270}=q_-^{270}-q_+^{270}.
\]
Les 324 parcours orientés sont reliés bijectivement à leurs cinq coordonnées historiques par l'identité canonique de parcours. Cette liaison utilise uniquement la structure généalogique et ne consulte aucune grandeur métrologique.

\Needspace{7\baselineskip}
\subsection{Angle direct, contre-angle et canaux conjugués}
La coordonnée angulaire directe procède de la constante de structure fine déjà construite :

\[
 A=1000\alpha_{\rm HMT}
 =7.297352569283800997285105472380663\ldots{}^\circ.
\]
Le contre-angle ne s'obtient ni à partir d'une masse ni à partir de la fonctionnelle de Barbero--Immirzi. Soit

\[
 D_A=\exp\!\left(-\frac{100\pi}{9}\alpha_{\rm HMT}\right),
\]
\[
 \mathcal C_\pi(\alpha_{\rm HMT})
 =169\alpha_{\rm HMT}^{6}
 +\frac{D_A\alpha_{\rm HMT}^{7}}
        {1-D_A\alpha_{\rm HMT}},
\]
\[
 y_*=\frac{\pi}{90}(H_5+\alpha_{\rm HMT})
      -\mathcal C_\pi(\alpha_{\rm HMT}),
 \qquad
 C^*=\frac{180}{\pi}y_*.
\]
Par conséquent,

\[
 C^*=2.123738338968645724261951380393\ldots{}^\circ.
\]
La coordonnée \(A\) est le mode angulaire commun, et \(C^*\) le mode différentiel orienté. Leur passage à la carte archimédienne est

\[
 x=A_{\rm rad}=\frac{\pi A}{180},
 \qquad
 y=C^*_{\rm rad}=\frac{\pi C^*}{180}.
\]
Dans l'algèbre réelle engendrée par l'involution \(R=\operatorname{diag}(1,-1)\), avec \(P_\pm=(I\pm R)/2\), le bloc biangulaire est

\[
 B_{A,C^*}=AI+C^*R
 =(A+C^*)P_++(A-C^*)P_-.
\]
L'angle plus le contre-angle et l'angle moins le contre-angle sont ainsi les deux valeurs propres d'un même objet. L'exponentiation produit les canaux

\[
 q_+=e^{-x-y},\qquad q_-=e^{-x+y}.
\]
La carte est inversible :

\[
 x=-\frac12\log(q_+q_-),
 \qquad
 y=\frac12\log\frac{q_-}{q_+}.
\]
La transformation \(C^*\mapsto-C^*\) conserve les projecteurs et échange les canaux. Après la construction régionale, la coordonnée de retour satisfait

\[
 \eta_{\rm ret}=\frac{C^*}{2}-\alpha_{\rm HMT}\,{\rm deg},
\]
sans constituer une seconde règle de sélection de \(C^*\).

\Needspace{7\baselineskip}
\subsection{Caractère central}
Avec

\[
 \Theta_0=
 \left(A_{\rm rad},\frac{C^*_{\rm rad}}6,
 \Delta_4,s_{120},\zeta_{270}\right),
\]
le caractère central est

\[
 \chi_0(\sigma)=
 \exp\!\left(
 n_AA_{\rm rad}+\frac{s_C}{6}C^*_{\rm rad}
 +k_{\Delta_4}\Delta_4+b_{120}s_{120}+b_\zeta\zeta_{270}
 \right).
\]
Par additivité de la signature,

\[
 \chi_0(\sigma_1+\sigma_2)=\chi_0(\sigma_1)\chi_0(\sigma_2).
\]
L'unité physique n'intervient pas dans cette exponentielle : \(\chi_0\) est un automorphisme positif adimensionnel de la droite de masse.

\Needspace{7\baselineskip}
\section{Hiérarchie harmonique globale des tours}
Avec les coordonnées (x=\allowbreak{}A\_\{\textbackslash{}rm rad\}) et (y=\allowbreak{}C\textasciicircum{}*\_\{\textbackslash{}rm rad\}) déjà construites, soient

\[
 q_\pm=\exp\!\left[-\frac{\pi}{180}(A\pm C^*)\right]
 =e^{-x\mp y}.
\]
Dans une convention globale cohérente,

\[
 T=e^{-xI-yR},\qquad
 c_n=\operatorname{Tr}(T^n)=q_-^n+q_+^n,
\]
\[
 s_n=-\operatorname{Tr}(RT^n)=q_-^n-q_+^n.
\]
De manière équivalente, sous forme opératorielle,

\[
 c_n=2e^{-nx}\cosh(ny),\qquad
 s_n=2e^{-nx}\sinh(ny).
\]
L'involution \(C^*\mapsto-C^*\) fixe les projecteurs et permute les coefficients spectraux ou, de manière équivalente, les canaux associés. On en déduit

\[
 c_{m+n}=\frac{c_mc_n+s_ms_n}{2},\qquad
 s_{m+n}=\frac{s_mc_n+c_ms_n}{2},
\]
\[
 c_n^2-s_n^2=4e^{-2nx},
\]
et la récurrence commune d'ordre deux. Les hauteurs admissibles forment

\[
 \mathfrak S=\langle90,120\rangle
 =\{90,120\}\cup\{30r:r\ge6\}.
\]
Les valeurs \(90,120,180,210,240,270,360,420\) sont des témoins généalogiques, et non la tour complète. Avant le quotient, la hauteur 360 conserve deux histoires :

\[
 k[\mathfrak S_0]
 \simeq k[X_{90},Y_{120}]/(X_{90}^4-Y_{120}^3).
\]
Le relèvement général s'écrit

\[
 \chi_{\rm full}(\sigma,\eta)
 =\chi_0(\sigma)
  \exp\!\left(\sum_{h\in\mathfrak S}\eta_hs_h\right),
\]
avec

\[
 N_{120}=b_{120}+\eta_{120}.
\]
La somme numérique n'efface pas la double provenance : \(b_{120}\) provient de l'enregistrement et \(\eta_{120}\) du raffinement spectral. Elle ne peut pas non plus absorber \(b_\zeta\zeta_{270}\) dans \(\eta_{270}s_{270}\).

\Needspace{7\baselineskip}
\section{Branche absolue : action, horloge et droite de masse}
Le lecteur déterminantal produit

\[
 \operatorname{SNF}(H_4)=\operatorname{diag}(1,2,2,4),\qquad
 |\operatorname{coker}H_4|=16,
\]
\[
 \Sigma_H=\varphi\alpha_{\rm HMT}^{16},\qquad
 \operatorname{ord}_{10}(\Sigma_H)=-34.
\]
Les deux sections orientées d'une seule droite d'action sont

\[
 \hbar_{\rm pre}=H_5\,10^{-34}\mathcal U_S,
 \qquad
 \hbar_{\rm ret}=\eta_{\rm ret}\,10^{-34}\mathcal U_S.
\]
Leur carte additive fine et leur carte multiplicative sont

\[
 \delta_{2w}=H_5-\eta_{\rm ret},\qquad
 R_{\rm act}=\frac{H_5}{\eta_{\rm ret}}>0.
\]
Il ne s'agit pas de deux constantes de Planck. La décennie commune fixe l'ordre absolu ; le quotient transporte l'information relative d'aller/retour.

Une fois \(t_0\in\mathcal L_T^+\) déclaré, CT--108 fixe

\[
 \omega_0=\frac{2\pi}{108t_0},\qquad
 \varepsilon_{\rm clk}=\hbar_{\rm ret}\omega_0,
 \qquad
 m_{\rm clk}=\varepsilon_{\rm clk}c_{\rm int}^{-2}.
\]
Le caractère agit par homothétie sur \(\mathcal L_M\) ; il ne devient pas \(\hbar\). Dans les quotients de masses, tout facteur commun se simplifie :

\[
 \frac{m_Y}{m_X}
 =\chi_{\rm full}(\sigma_Y-\sigma_X,\eta_Y-\eta_X)
 R_{\rm act}^{\kappa_Y-\kappa_X}.
\]
Ainsi, \(10^{-34}\) fixe l'échelle absolue, mais ne corrige pas un résidu relatif. Une modification relative ne peut provenir que du parcours, de l'horloge spécifique, du lecteur ou des tours.

Une distinction reste nécessaire. HMT construit une section interne une fois \(t_0\) déclaré ; la coordonnée MeV exige une carte. Pour déterminer complètement la carte absolue sans convention externe, il faut construire

\[
 \Gamma\longmapsto t_\Gamma
 \quad\text{ou}\quad
 \Gamma\longmapsto\omega_\Gamma.
\]
La section 10 définit un lecteur d'horloge fondé sur les retours et l'holonomie. Sa normalisation absolue est déterminée lorsque l'opérateur de parcours fixe une section positive de la droite temporelle.

\Needspace{7\baselineskip}

\endgroup

La forme normale, le conoyau d’ordre seize, les ablations et la distinction entre \(h\) et \(\hbar\) sont démontrés intégralement dans \cref{chap:escala-accion}. Ce résultat déjà construit est utilisé ici pour fixer la branche absolue de la loi, sans reprendre sa dérivation.
\begingroup
\let\chapter\section
\let\section\subsection
\let\subsection\subsubsection
\let\subsubsection\paragraph
\let\clearpage\relax
\chapter{Coordonnées angulaires conjuguées et décomposition spectrale}
\label{chap:canales}

Le passage de la clôture arithmétique de \(\alpha\) à la hiérarchie
dimensionnelle s'effectue au moyen de deux coordonnées angulaires. La première provient
directement de \(\alphaHMT\) ; la seconde est la phase orientée produite par
la réalisation analytique régionale de la microfibre quintuple. Toutes deux arrivent en
mécanique dimensionnelle déjà construites, avant l'introduction de la carte exponentielle
qui engendrera les moments globaux.

\section{Transfert des coordonnées angulaires dérivées dans HMT}

Les Parties I--III ont construit des dépendances distinctes ayant une racine commune dans
\(\alpha_{\rm HMT}\). Le caractère pré-retour provient de
\[
(\alpha_{\rm HMT},\varphi;\,9,5,4,12,54,\text{ règle graduée})
\longrightarrow H_5,
\]
tandis que la coordonnée angulaire et le déterminant proviennent de
\[
\alpha_{\rm HMT}\longrightarrow
A=1000\alpha_{\rm HMT}\,{}^\circ,
\qquad
(\alpha_{\rm HMT},\pi)\longrightarrow
D_A=e^{-100\pi\alpha_{\rm HMT}/9}.
\]
La microfibre régionale fixe, par une dépendance distincte,
\(\mathcal F_\pi\to\rho_\pi=169\), et avec elle
\[
\mathcal C_\pi(\alpha_{\rm HMT})
=169\alpha_{\rm HMT}^{6}
+\frac{D_A\alpha_{\rm HMT}^{7}}
       {1-D_A\alpha_{\rm HMT}}.
\]
La branche de \(H_5\) et la correction régionale ne confluent qu'à l'étape
suivante :
\[
\begin{aligned}
y_*&=\frac{\pi}{90}(H_5+\alpha_{\rm HMT})
     -\mathcal C_\pi(\alpha_{\rm HMT}),\\
C^*&=\frac{180}{\pi}y_*.
\end{aligned}
\]
Désormais, \(A\) désigne la coordonnée angulaire directe et \(C^*\) la
coordonnée angulaire régionale conjuguée ; ces dénominations fixent leurs types
sans introduire une seconde règle de sélection.
Par conséquent, \(H_5\) n'est pas une sortie de \(A\). Ici, \(y_*\) est la coordonnée
analytique en radians et \(C^*\) sa coordonnée sexagésimale. La mécanique
dimensionnelle reçoit ces sorties déjà déterminées ; elle ne les redéfinit pas.
La coordonnée \(A\) n'est pas une entrée directe de la formule de \(C^*\) : les deux
objets sont couplés pour la première fois dans \((A,C^*)\mapsto q_\pm\). L'égalité
\(\alpha_{\rm HMT}=A/1000\) n'est que le changement de coordonnée inverse.
Dans cette Partie, nous conservons les notations
\[
\begin{aligned}
 A&=7.297352569283800997285105472380663^\circ,\\
 \Cstar&=2.123738338968645724261951380393\ldots{}^\circ.
\end{aligned}
\]
L'arrondi utilisé dans les tableaux est
\(2.123738338968646^\circ\). L'entier \(169\), le déterminant \(D_A\) et
la récurrence régionale appartiennent à la démonstration précédente ; aucune
comparaison SI ni la fonctionnelle hexade--octade n'interviennent dans leur génération.

L'identité
angulaire
\[
 \etaRet=\frac{\Cstar}{2}-\alphaAngle,
 \qquad
 \alphaAngle:=\alphaHMT\,{\rm deg},
\]
définit la coordonnée angulaire adimensionnelle postérieure au retour ; elle n'est pas utilisée
pour engendrer \(\Cstar\). De manière équivalente,
une fois la construction régionale effectuée,
\[
 \Cstar=2(\etaRet+\alphaAngle).
\]
Cette dernière égalité est la forme inverse de la même relation et non une
seconde règle de sélection de \(C^*\).

La réalisation régionale corrige \(H_5\) avant de produire \(\etaRet\) ; par
conséquent, \(H_5\),
\(\etaRet\) et \(\Cstar\) ont des types mathématiques distincts. Le symbole
\(\hbarHMT\) est réservé à la grandeur sur la droite d'action obtenue
seulement après application de l'échelle \(10^{-34}\mathcal U_S\). La comparaison
avec une mantisse du SI appartient au contrôle métrologique externe et laisse
invariante la construction précédente.

\section{Phases internes et carte exponentielle}

Les phases abstraites associées aux deux angles sont
\[
 a=\frac A{360},\qquad c=\frac{\Cstar}{360}
 \quad\text{dans }\RR/\ZZ.
\]
Avant d'appliquer une exponentielle réelle, on choisit les représentants canoniques
\(\widetilde a,\widetilde c\in[0,1)\). L'exponentielle agit sur ces
représentants, non directement sur les classes du quotient.
La conversion en radians n'est introduite que lors du choix de la carte
archimédienne
\[
 x:=A_{\rm rad}=\frac{\pi A}{180},\qquad
 y:=\Cstarrad=\frac{\pi\Cstar}{180},
\]
\[
 q_-=e^{-x+y},\qquad q_+=e^{-x-y}.
\]
Le facteur \(\pi/180\) appartient à ce changement analytique de coordonnées. La
clôture arithmétique de \(\alpha\), formulée avant cette conversion, est
indépendante de ce choix de carte angulaire.

\section{Algèbre involutive centrale}

Considérons la représentation régulière bidimensionnelle de l'algèbre réelle
engendrée par une involution. Dans une base spectrale, nous fixons
\[
 R=\operatorname{diag}(1,-1),\qquad I=I_2,
\]
de sorte que les deux sous-espaces propres ont une multiplicité un. Soient
\[
 P_+=\frac{I+R}{2},\qquad P_-=\frac{I-R}{2}
\]
ses projecteurs spectraux. L'élément central déterminé par les
coordonnées angulaires est
\[
 B_{A,\Cstar}=AI+\Cstar R
 =(A+\Cstar)P_++(A-\Cstar)P_-.
\]
L'involution \(\Cstar\mapsto-\Cstar\) conserve les projecteurs fixes
\(P_\pm\) et échange les deux coefficients spectraux ---de manière équivalente,
les canaux associés---. En appliquant l'exponentielle, on obtient
\[
 T=e^{-xI-yR}=q_+P_++q_-P_-.
\]
Par conséquent, les deux canaux sont les valeurs propres d'un unique élément de
l'algèbre engendrée par \(I\) et \(R\).

\section{Inversibilité de la carte des canaux}

\begin{theorem}[Reconstruction des coordonnées angulaires]
Para \(q_\pm>0\),
\[
 x=-\frac12\log(q_+q_-),
 \qquad
 y=\frac12\log\frac{q_-}{q_+}.
\]
Par conséquent, l'application
\((A,\Cstar)\mapsto(q_+,q_-)\) est injective. L'échange
\(q_+\leftrightarrow q_-\) conserve \(A\) et transforme
\(\Cstar\mapsto-\Cstar\).
\end{theorem}
\begin{proof}
La multiplication et le quotient des deux exponentielles donnent
\(q_+q_-=e^{-2x}\) et \(q_-/q_+=e^{2y}\). En prenant les logarithmes, on retrouve
\(x\) et \(y\). Comme \(\pi/180\neq0\), les deux angles sont déterminés.
\end{proof}

\section{Identités d'addition et de récurrence}

Dans cette représentation bidimensionnelle, les moments pair et impair de l'opérateur
\(T\) sont
\[
 c_n=\Tr(T^n)=q_-^n+q_+^n,
 \qquad
 s_n=-\Tr(RT^n)=q_-^n-q_+^n.
\]

L'hypothèse de multiplicité un fait partie du domaine de ces formules.
Dans une représentation de multiplicités \(m_+\) et \(m_-\), les traces
correspondantes seraient
\(m_+q_+^n+m_-q_-^n\) et
\(m_-q_-^n-m_+q_+^n\), respectivement.

\begin{theorem}[Lois de composition des moments]
Pour \(m,n\geq0\),
\[
 c_{m+n}=\frac{c_mc_n+s_ms_n}{2},
 \qquad
 s_{m+n}=\frac{s_mc_n+c_ms_n}{2},
\]
\[
 c_n^2-s_n^2=4e^{-2nx},
\]
et les deux suites satisfont
\[
 X_{n+2}=c_1X_{n+1}-e^{-2x}X_n.
\]
De plus,
\[
 \partial_Ac_n=-\frac{n\pi}{180}c_n,
 \quad
 \partial_As_n=-\frac{n\pi}{180}s_n,
\]
\[
 \partial_{\Cstar}c_n=\frac{n\pi}{180}s_n,
 \quad
 \partial_{\Cstar}s_n=\frac{n\pi}{180}c_n.
\]
\end{theorem}
\begin{proof}
On développe les produits de
\(q_-^m\pm q_+^m\) avec \(q_-^n\pm q_+^n\). L'identité quadratique résulte
de \((u+v)^2-(u-v)^2=4uv\). La récurrence correspond au polynôme
caractéristique \(z^2-c_1z+q_-q_+\), et les dérivées s'obtiennent par la
règle de la chaîne.
\end{proof}

\begin{corollary}[Détermination de la hiérarchie complète]
Une fois \(A\) et \(\Cstar\) fixés, tous les couples \((c_n,s_n)\) sont
déterminés par la récurrence. Aucune espèce physique n'introduit de canal
supplémentaire.
\end{corollary}

Enfin,
\[
 s_n=2e^{-nx}\sinh(ny),
 \qquad
 c_n=2e^{-nx}\cosh(ny).
\]
La coordonnée \(A\) détermine la décroissance commune et \(\Cstar\) la séparation
orientée. L'ordre décimal de la droite d'action n'est pas introduit dans ces
formules. Il se déduit dans le \cref{chap:escala-accion} à partir du conoyau
entier du bloc local de Hadamard. Le choix ultérieur d'une base
exprimée en \(\mathrm{J\,s}\) appartient au transducteur métrologique et ne
modifie pas les identités adimensionnelles.

\endgroup
\begingroup
\let\chapter\section
\let\section\subsection
\let\subsection\subsubsection
\let\subsubsection\paragraph
\let\clearpage\relax
\section{Ellipse d'action, radian et double retour d'holonomie}
\label{sec:elipse-accion-holonomia-rev8}
\index{action!ellipse}
\index{radian!secteur d'action}
\index{Möbius!double retour}

Le théorème focal précédent fixe la forme angulaire et le lecteur déterminantiel
fixe l'échelle d'action. Définissons
\[
 \eta=\operatorname{artanh}\frac{C^*}{A},
 \qquad
 a_S=\sqrt{2\hbar_{\rm HMT}^{\rm ret}e^\eta},
 \qquad
 b_S=\sqrt{2\hbar_{\rm HMT}^{\rm ret}e^{-\eta}}.
\]
L'orbite symplectique
\[
 Q(\theta)=a_S\cos\theta,
 \qquad
 P(\theta)=b_S\sin\theta
\]
conserve simultanément le rapport de forme du bloc angulaire et l'aire
fixée par la droite d'action.

\begin{theorem}[Aire balayée par phase]
\label{thm:area-fase-elipse-rev8}
Pour tout \(\theta\),
\[
 \operatorname{Area}(\theta)
 =\frac12\int_0^\theta(QP'-PQ')\,dt
 =\hbar_{\rm HMT}^{\rm ret}\theta.
\]
Par conséquent,
\[
 \operatorname{Area}(1\ {\rm rad})=\hbar_{\rm HMT}^{\rm ret},
 \qquad
 \operatorname{Area}(2\pi)=h_{\rm HMT}^{\rm ret}.
\]
\end{theorem}

\begin{proof}
Comme \(QP'-PQ'=a_Sb_S\) et
\(a_Sb_S=2\hbar_{\rm HMT}^{\rm ret}\), l'intégrale vaut
\(\hbar_{\rm HMT}^{\rm ret}\theta\). La seconde identité utilise
\(h_{\rm HMT}^{\rm ret}=2\pi\hbar_{\rm HMT}^{\rm ret}\).
\end{proof}

APP fournit l'aire orientée ; \(A,C^*\) fixent la forme et les foyers ; la branche
déterminantielle fixe l'échelle. Le tour primitif enferme \(h\) et un radian
de phase balaie \(\hbar\). Le radian ordinaire reste
\(1\,\mathrm{rad}=180^\circ/\pi\) ; la construction ajoute sa généalogie comme
secteur unitaire d'action, non une autre constante angulaire.

Pour une section persistante d'holonomie
\(\varepsilon_\gamma\in\{+1,-1\}\), la condition
\[
 \varepsilon_\gamma e^{iJ_\gamma/\hbar}=1
\]
distingue deux clôtures. Dans le cylindre, \(\varepsilon_\gamma=+1\) et le cycle
primitif transporte \(J=2\pi\hbar=h\). Dans le ruban de Möbius,
\(\varepsilon_\gamma=-1\) : un tour visible change de feuillet et transporte
\(J=\pi\hbar=h/2\) ; le second rétablit l'orientation et complète \(h\). Par
radian de la phase interne, on transporte \(\hbar\) dans les deux cas ; le facteur
\(1/2\) appartient à la projection sur le tour visible de la base de Möbius.

\begin{figure}[htbp]
\centering
\resizebox{.98\textwidth}{!}{\begin{tikzpicture}[
  font=\scriptsize,
  box/.style={draw,rounded corners=2mm,align=center,minimum height=.9cm,
    inner xsep=5pt},
  arr/.style={-{Stealth[length=2.2mm]},thick}]

  \begin{scope}[shift={(0,0)}]
    \draw[->,hmtgray] (-3.0,0)--(3.0,0) node[right]{\(u_+\)};
    \draw[->,hmtgray] (0,-2.0)--(0,2.0) node[above]{\(u_-\)};
    \draw[hmtblue,very thick] (0,0) ellipse (2.65cm and 1.45cm);
    \fill[hmtgold] (-1.85,0) circle (2pt) node[below]{\(F_-^{(\lambda)}\)};
    \fill[hmtgold] (1.85,0) circle (2pt) node[below]{\(F_+^{(\lambda)}\)};
    \node[text=hmtblue,font=\bfseries] at (0,2.55) {ellipse des taux};
    \node[align=center,text=hmtgray] at (0,-2.4)
      {\(q_+=e^{-a_\lambda^2}\), \(q_-=e^{-b_\lambda^2}\)};
  \end{scope}

  \draw[arr,hmtgold] (3.3,0)--node[above,align=center]{même forme\\échelle d’action} (6.0,0);

  \begin{scope}[shift={(9.0,0)}]
    \draw[->,hmtgray] (-2.8,0)--(2.8,0) node[right]{\(Q\)};
    \draw[->,hmtgray] (0,-2.0)--(0,2.0) node[above]{\(P\)};
    \fill[hmtlightgold] (0,0)--(2.5,0)
      arc[start angle=0,end angle=57.2958,x radius=2.5cm,y radius=1.35cm]--cycle;
    \draw[hmtviolet,very thick] (0,0) ellipse (2.5cm and 1.35cm);
    \draw[hmtgold,thick] (0,0)--(2.5,0);
    \draw[hmtgold,thick] (0,0)--({2.5*cos(57.2958)},{1.35*sin(57.2958)});
    \node[text=hmtviolet,font=\bfseries] at (0,2.55) {ellipse d’action};
    \node[align=center] at (.7,.48) {\(\theta=1\)\\aire \(\hbar\)};
    \node[align=center,text=hmtgray] at (0,-2.4) {tour \(2\pi\) : aire \(h\)};
  \end{scope}

  \node[box,draw=hmtblue,fill=hmtlightblue,minimum width=3.2cm] (cyl) at (3.2,-4.3)
    {cilindro \(\varepsilon=+1\)\\\(2\pi\mapsto h\)};
  \node[box,draw=hmtviolet,fill=hmtlightviolet,minimum width=4.2cm] (mob) at (9.0,-4.3)
    {Möbius \(\varepsilon=-1\)\\\(2\pi\mapsto h/2\), \(4\pi\mapsto h\)};
  \draw[arr,hmtgray] (cyl)--node[
    above,align=center,text width=1.9cm,fill=white,inner sep=1pt
  ]{changement d’\\holonomie} (mob);
  \node[align=center,text=hmtgray,font=\footnotesize] at (6.1,-5.55)
    {La phase interne transporte \(\hbar\) par radian ;\\le facteur deux appartient au tour visible.};
\end{tikzpicture}
}
\caption{L'ellipse des taux fixe la forme ; son relèvement à la droite d'action fixe
l'aire. La clôture cylindrique et le double retour de Möbius distinguent la phase
interne du tour visible.}
\label{fig:elipse-accion-holonomia-rev8}
\end{figure}

\subsection{Géométrie de Hilbert--Beltrami--Klein dans l'ellipse d'action}
\label{sec:hilbert-beltrami-klein-elipse-accion}

La même ellipse qui porte l'aire symplectique admet une géométrie projective
intrinsèque. Soit
\[
 \mathbb E_{\rm act}
 :=\left\{(Q,P)\in\mathbb R^2:
 \frac{Q^2}{a_S^2}+\frac{P^2}{b_S^2}<1\right\},
 \qquad
 \mathcal L_{\rm act}(Q,P)
 :=\left(\frac{Q}{a_S},\frac{P}{b_S}\right).
\]
L'application linéaire \(\mathcal L_{\rm act}\) identifie
\(\mathbb E_{\rm act}\) au disque unité \(\mathbb D\). Pour deux points
distincts \(x,y\in\mathbb E_{\rm act}\), soient \(a,b\in\partial
\mathbb E_{\rm act}\) les intersections de la droite \(xy\) avec la frontière,
ordonnées comme \(a,x,y,b\). Nous définissons
\begin{equation}
 d_{\rm H}^{\rm act}(x,y)
 :=\frac12\log
 \frac{\lVert y-a\rVert\,\lVert x-b\rVert}
      {\lVert x-a\rVert\,\lVert y-b\rVert}.
 \label{eq:distancia-hilbert-elipse-accion}
\end{equation}
L'expression est le birapport des quatre points colinéaires et est donc
indépendante de la norme auxiliaire utilisée sur cette droite.

\begin{proposition}[Réalisation de Beltrami--Klein de l'ellipse d'action]
\label{prop:beltrami-klein-elipse-accion}
Le couple \((\mathbb E_{\rm act},d_{\rm H}^{\rm act})\) est isométrique, par
\(\mathcal L_{\rm act}\), au modèle de Beltrami--Klein du plan hyperbolique.
Ses géodésiques non orientées sont exactement les cordes ouvertes de
l'ellipse. La frontière \(\partial\mathbb E_{\rm act}\) conserve, de manière
simultanée et typée, l'orbite symplectique
\[
 \gamma_{\rm act}(\theta)
 =(a_S\cos\theta,b_S\sin\theta),
 \qquad
 \frac12\int_0^\theta(QP'-PQ')\,dt
 =\hbar_{\rm HMT}^{\rm ret}\theta.
\]
\end{proposition}

\begin{proof}
Les transformations projectives préservent les birapports. Comme
\(\mathcal L_{\rm act}\) est une transformation linéaire inversible qui envoie
l'ellipse sur le disque unité, elle transporte
\eqref{eq:distancia-hilbert-elipse-accion} vers la distance de Hilbert du
disque. Celle-ci est la métrique de Beltrami--Klein, dont les géodésiques sont les cordes.
L'identité symplectique de la frontière est le
\cref{thm:area-fase-elipse-rev8} écrit dans la carte normalisée.
\end{proof}

La carte complète est donc typée par
\[
 \mathfrak R_{\rm act}:
 (A,C^*,\hbar_{\rm HMT}^{\rm ret})
 \longmapsto
 (\mathbb E_{\rm act},d_{\rm H}^{\rm act},
  dP\wedge dQ,\gamma_{\rm act}).
\]
La métrique projective gouverne l'intérieur ; la forme symplectique et la loi
d'aire gouvernent la frontière parcourue. Une réalisation comme région physique ou
comme domaine de matrices de covariance se formule par un foncteur ultérieur
qui conserve explicitement ces deux structures ; la construction précédente
fixe déjà son domaine mathématique et les invariants que ce foncteur doit préserver.

\endgroup
\begingroup
\let\chapter\section
\let\section\subsection
\let\subsection\subsubsection
\let\subsubsection\paragraph
\let\clearpage\relax
\subsection[Identité autoduale de Planck sur la périodicité d'ordre 108]{Identité autoduale de
Planck dans la réalisation circulaire de l'endomorphisme temporel d'ordre 108}
\label{sec:identidad-autodual-planck-ct108-rev3}
\index{Planck!identité autoduale}
\index{calendrier d’ordre 108!réalisation circulaire}
\index{Compton!lectures réduite et complète}

La droite d’action construite dans la section précédente admet une
réalisation circulaire compatible avec la périodicité temporelle d’ordre 108.
Cette réalisation coordonne trois objets qui possèdent déjà leurs types propres : la
section d’action, la géométrie d’un tour et la transduction
gravitationnelle. La coordination qui en résulte fournit une identité de
Planck interne et prépare la droite positive sur laquelle agira ensuite
l’opérateur général de masse.

\subsubsection{Cartes orientées d’action et réalisation circulaire déclarée}

Soient
\[
 \mathcal L_S^+,\quad \mathcal L_T^+,\quad \mathcal L_C^+,
 \quad \mathcal L_L^+,\quad \mathcal L_M^+,\quad \mathcal L_G^+
\]
les demi-droites positives d’action, de temps, de vitesse, de longueur et de masse.
Les coordonnées antérieure et postérieure au retour s’écrivent
\begin{equation}
 \hbar_{\rm pre}=H_5\,10^{-34}\mathcal U_S,
 \qquad
 \hbar_{\rm ret}=\eta_{\rm ret}\,10^{-34}\mathcal U_S,
 \qquad
 R_{\rm act}=\frac{\hbar_{\rm pre}}{\hbar_{\rm ret}}
 =\frac{H_5}{\eta_{\rm ret}}.
 \label{eq:planck-ct108-cartas-pre-ret-rev3}
\end{equation}
L’indice \(a\in\{\mathrm{pre},\mathrm{ret}\}\) désignera l’une quelconque
des deux cartes. Dans chacune, la distinction exacte est conservée
\begin{equation}
 h_a=2\pi\hbar_a .
 \label{eq:planck-ct108-h-hbar-rev3}
\end{equation}

Fixons une durée élémentaire \(t_0\in\mathcal L_T^+\), une vitesse
interne \(c_{\rm int}\in\mathcal L_C^+\) et
\(\ell_0=c_{\rm int}t_0\). La réalisation circulaire déclarée de la
périodicité temporelle d’ordre 108 est déterminée par
\begin{equation}
 T_{108}=108t_0,
 \qquad
 \omega_{108}=\frac{2\pi}{T_{108}},
 \qquad
 R_{108}=\frac{c_{\rm int}}{\omega_{108}}
 =\frac{54}{\pi}\ell_0,
 \label{eq:planck-ct108-realizacion-circular-rev3}
\end{equation}
et par la longueur du tour complet
\begin{equation}
 C_{108}=2\pi R_{108}=c_{\rm int}T_{108}=108\ell_0.
 \label{eq:planck-ct108-perimetro-rev3}
\end{equation}
Ainsi, \(R_{108}\) est le rayon de la réalisation et \(C_{108}\) son
périmètre. La fréquence angulaire, la fréquence cyclique et les deux
constantes d’action sont reliées par
\(\omega_{108}=2\pi/T_{108}\) et \(h_a=2\pi\hbar_a\).

\subsubsection{Correspondance entre la périodicité d’ordre 108 et les lectures de Compton}

Dans la carte \(a\), le quantum circulaire et sa section de masse sont
\begin{equation}
 E_{108,a}=\hbar_a\omega_{108},
 \qquad
 m_{108,a}=\frac{E_{108,a}}{c_{\rm int}^{2}}
 =\frac{\hbar_a\omega_{108}}{c_{\rm int}^{2}}.
 \label{eq:planck-ct108-cuanto-circular-rev3}
\end{equation}

\begin{theorem}[Identité entre la périodicité circulaire et les lectures de Compton]
\label{thm:planck-ct108-compton-rev3}
Pour chaque carte \(a\in\{\mathrm{pre},\mathrm{ret}\}\), les lectures de
Compton réduite et complète satisfont
\begin{equation}
 \boxed{
 \bar\lambda_{\rm C}(m_{108,a})
 =\frac{\hbar_a}{m_{108,a}c_{\rm int}}=R_{108},
 \qquad
 \lambda_{\rm C}(m_{108,a})
 =\frac{h_a}{m_{108,a}c_{\rm int}}=C_{108}.}
 \label{eq:planck-ct108-identidad-compton-rev3}
\end{equation}
\end{theorem}

\begin{proof}
La substitution de
\eqref{eq:planck-ct108-cuanto-circular-rev3} produit
\[
 \frac{\hbar_a}{m_{108,a}c_{\rm int}}
 =\frac{c_{\rm int}}{\omega_{108}}=R_{108}.
\]
La deuxième égalité résulte de l’application simultanée de
\eqref{eq:planck-ct108-h-hbar-rev3} et de
\eqref{eq:planck-ct108-perimetro-rev3}.
\end{proof}

Le théorème est homogène en
\((\hbar_a,c_{\rm int},t_0)\) : la même périodicité détermine le rayon,
le périmètre et les deux cartes de Compton. La réalisation circulaire est la
primitive géométrique déclarée ; les égalités précédentes sont ses
conséquences exactes.

\subsubsection{Sélecteur circulaire et unité interne de Planck}

La droite gravitationnelle est coordonnée à la carte \(a\) par le transducteur
homogène
\begin{equation}
 G_a(\theta)=\theta\,
 \frac{c_{\rm int}^{5}t_0^2}{\hbar_a}
 \in\mathcal L_G^+,
 \qquad \theta\in\mathbb R_{>0}.
 \label{eq:planck-ct108-transductor-rev3}
\end{equation}
Son rayon gravitationnel réduit, évalué sur le quantum circulaire, est
\begin{equation}
 r_{{\rm g},a}(\theta)
 :=\frac{G_a(\theta)m_{108,a}}{c_{\rm int}^{2}}
 =\theta\frac{\pi}{54}\ell_0.
 \label{eq:planck-ct108-radio-gravitatorio-rev3}
\end{equation}

\begin{proposition}[Sélecteur circulaire de Planck]
\label{prop:selector-circular-planck-ct108-rev3}
Au sein de la réalisation circulaire
\eqref{eq:planck-ct108-realizacion-circular-rev3}, la condition
\(r_{{\rm g},a}(\theta)=R_{108}\) sélectionne une unique coordonnée
positive,
\begin{equation}
 \boxed{\theta_{108}^{\rm circ}=\left(\frac{54}{\pi}\right)^2.}
 \label{eq:theta-circular-planck-ct108-rev3}
\end{equation}
La coordonnée est indépendante de la carte orientée \(a\).
\end{proposition}

\begin{proof}
Les équations
\eqref{eq:planck-ct108-realizacion-circular-rev3} et
\eqref{eq:planck-ct108-radio-gravitatorio-rev3} réduisent la condition à
\[
 \theta\frac{\pi}{54}\ell_0=\frac{54}{\pi}\ell_0.
\]
La positivité de \(\ell_0\) et la stricte monotonie de l’application
\(\theta\mapsto r_{{\rm g},a}(\theta)\) fournissent la solution et son
unicité.
\end{proof}

Soit \(G_{108,a}^{\rm circ}:=G_a(\theta_{108}^{\rm circ})\). Les unités internes
associées à cette carte satisfont
\begin{align}
 \ell_{{\rm P},a}^{\rm circ}
 &:=\sqrt{\frac{\hbar_a G_{108,a}^{\circ}}{c_{\rm int}^{3}}}
 =R_{108},
 &
 t_{{\rm P},a}^{\circ}
 &:=\sqrt{\frac{\hbar_a G_{108,a}^{\circ}}{c_{\rm int}^{5}}}
 =\omega_{108}^{-1},
 \label{eq:planck-ct108-longitud-tiempo-rev3}\\
 m_{{\rm P},a}^{\circ}
 &:=\sqrt{\frac{\hbar_a c_{\rm int}}{G_{108,a}^{\circ}}}
 =m_{108,a},
 &
 E_{{\rm P},a}^{\circ}
 &:=m_{{\rm P},a}^{\circ}c_{\rm int}^{2}=E_{108,a}.
 \label{eq:planck-ct108-masa-energia-rev3}
\end{align}
L’égalité des rayons est donc publiée comme
\begin{equation}
 R_{108}=\bar\lambda_{\rm C}=r_{\rm g}=\ell_{\rm P}^{\rm circ},
 \qquad
 C_{108}=\lambda_{\rm C}=2\pi r_{\rm g}=2\pi\ell_{\rm P}^{\rm circ}.
 \label{eq:planck-ct108-tres-lectores-rev3}
\end{equation}

La convention de ce théorème est le rayon gravitationnel réduit
\(r_{\rm g}=Gm/c_{\rm int}^{2}\). Le rayon de Schwarzschild satisfait
\(r_{\rm s}=2r_{\rm g}\) ; si cette seconde convention est adoptée, son croisement avec
\(\bar\lambda_{\rm C}\) se produit en
\(m=m_{{\rm P},a}^{\circ}/\sqrt2\). La déclaration du rayon fait donc
partie de la carte et évite de transférer par inadvertance un facteur deux
entre les deux identités.

La coordonnée \(\theta_{108}^{\rm circ}\) appartient à la spécialisation
circulaire déclarée. Le sélecteur gravitationnel général, formulé sur
l’holonomie enrichie et sa réponse spectrale, conserve sa résidence dans
la \cref{sec:selector-gravitatorio-rev11}. Ainsi, l’identité
circulaire fixe une normalisation interne de Planck et le critère général
fixe les conditions que doit satisfaire la réalisation physique de la
section gravitationnelle.

\subsubsection{Covariance pré-retour/post-retour}

Le produit des sections réciproques d’action et de gravité conserve la
géométrie circulaire :
\begin{equation}
 \frac{m_{108,\rm pre}}{m_{108,\rm ret}}=R_{\rm act},
 \qquad
 \frac{G_{108,\rm pre}^{\circ}}{G_{108,\rm ret}^{\circ}}
 =R_{\rm act}^{-1},
 \qquad
 \hbar_a G_{108,a}^{\circ}
 =\theta_{108}^{\rm circ}c_{\rm int}^{5}t_0^2.
 \label{eq:planck-ct108-covariancia-pre-ret-rev3}
\end{equation}
En particulier, \(\ell_{{\rm P},\rm pre}^{\rm circ}\) et
\(\ell_{{\rm P},\rm ret}^{\rm circ}\) coïncident avec \(R_{108}\). La mémoire
bidirectionnelle réside dans le quotient des masses et dans son réciproque
gravitationnel ; l’échelle géométrique réside dans le produit invariant.

La réciprocité admet une paramétrisation symétrique qui sépare l’échelle
commune de l’orientation du retour. Définissons
\begin{equation}
 \begin{aligned}
 \hbar_{\rm geom}
 &:=\sqrt{\hbar_{\rm pre}\hbar_{\rm ret}},
 &
 G_{108,\rm geom}^{\circ}
 &:=\sqrt{G_{108,\rm pre}^{\circ}G_{108,\rm ret}^{\circ}},
 &
 \xi_{\rm act}&:=\frac12\log R_{\rm act}.
 \end{aligned}
 \label{eq:planck-ct108-parametrizacion-simetrica-rev3}
\end{equation}
Alors
\[
 \hbar_{\rm pre}=\hbar_{\rm geom}e^{\xi_{\rm act}},
 \qquad
 \hbar_{\rm ret}=\hbar_{\rm geom}e^{-\xi_{\rm act}},
\]
tandis que les sections gravitationnelles conjuguées satisfont
\[
 G_{108,\rm pre}^{\circ}
 =G_{108,\rm geom}^{\circ}e^{-\xi_{\rm act}},
 \qquad
 G_{108,\rm ret}^{\circ}
 =G_{108,\rm geom}^{\circ}e^{\xi_{\rm act}}.
\]
La coordonnée \(\xi_{\rm act}\) conserve l’orientation
pré-retour/post-retour ; le produit \(\hbar_aG_{108,a}^{\circ}\)
conserve l’échelle géométrique. Cette paramétrisation des cartes d’action
n’introduit pas à elle seule des vitesses de propagation distinctes.

\subsubsection{Droite autoduale déployée des masses}

Une carte \(a\) étant fixée, soit l’algèbre déployée
\[
 \mathbb D_{\rm sp}:=\mathbb R[j]/(j^2-1),
 \qquad
 P_\pm:=\frac{I\pm j}{2},
\]
et considérons sa décomposition spectrale
\[
 \mathbb D_{\rm sp}
 =\mathbb R P_+\oplus\mathbb R P_-,
 \qquad
 P_\pm^2=P_\pm,\quad
 P_+P_-=0,\quad
 P_++P_-=I.
\]
La conjugaison déployée échange les projecteurs,
\[
 \overline{uP_++vP_-}=vP_++uP_-,
\]
et sa norme est
\[
 N_{\rm sp}(uP_++vP_-):=uv.
\]

Pour \(m\in\mathcal L_M^+\), définissons la section
\begin{equation}
 Z_a(m)
 :=\frac{m_{{\rm P},a}^{\circ}}{m}P_+
   +\frac{m}{m_{{\rm P},a}^{\circ}}P_- .
 \label{eq:planck-ct108-seccion-partida-masa-rev3}
\end{equation}
L’unité \(m_{{\rm P},a}^{\circ}\) détermine également l’involution
\begin{equation}
 \iota_{{\rm P},a}(m)
 =\frac{(m_{{\rm P},a}^{\circ})^2}{m}.
 \label{eq:planck-ct108-involucion-masas-rev3}
\end{equation}

\begin{theorem}[Droite autoduale déployée des masses]
\label{thm:planck-ct108-recta-partida-masa-rev3}
Pour tout \(m\in\mathcal L_M^+\),
\begin{equation}
 \boxed{
 N_{\rm sp}(Z_a(m))=1,
 \qquad
 \overline{Z_a(m)}
 =Z_a\!\left(\iota_{{\rm P},a}(m)\right).}
 \label{eq:planck-ct108-norma-conjugacion-rev3}
\end{equation}
La conjugaison possède un unique point fixe positif,
\[
 m=m_{{\rm P},a}^{\circ},
 \qquad
 Z_a(m_{{\rm P},a}^{\circ})=I.
\]
\end{theorem}

\begin{proof}
La norme est le produit des deux coordonnées réciproques :
\[
 \frac{m_{{\rm P},a}^{\circ}}m
 \frac m{m_{{\rm P},a}^{\circ}}=1.
\]
La conjugaison échange les deux coordonnées, ce qui équivaut à remplacer
\(m\) par \((m_{{\rm P},a}^{\circ})^2/m\). Le point fixe positif satisfait
\(m^2=(m_{{\rm P},a}^{\circ})^2\).
\end{proof}

Avec la rapidité
\begin{equation}
 x_a(m)=\log\frac{m}{m_{{\rm P},a}^{\circ}},
 \label{eq:planck-ct108-rapidez-masa-rev3}
\end{equation}
la section et sa conjuguée s’écrivent
\[
 Z_a(m)=e^{-x_a}P_++e^{x_a}P_-,
 \qquad
 \overline{Z_a(m)}
 =e^{x_a}P_++e^{-x_a}P_-.
\]
Leurs lecteurs pair et impair sont
\begin{equation}
 Q_a(m):=\frac{e^{x_a}+e^{-x_a}}2=\cosh x_a,
 \qquad
 A_a(m):=\frac{e^{x_a}-e^{-x_a}}2=\sinh x_a.
 \label{eq:planck-ct108-lectores-par-impar-rev3}
\end{equation}
Le premier conserve la distance au point autodual et le second conserve son
orientation.

Les lecteurs physiques réciproques associés à la même coordonnée sont
\[
 r_{{\rm g},a}(m)=\frac{G_{108,a}^{\circ}m}{c_{\rm int}^{2}},
 \qquad
 \bar\lambda_{{\rm C},a}(m)=\frac{\hbar_a}{mc_{\rm int}},
\]
et satisfont
\begin{equation}
 \frac{r_{{\rm g},a}(m)}{\bar\lambda_{{\rm C},a}(m)}
 =e^{2x_a(m)},
 \qquad
 \frac{\bar\lambda_{{\rm C},a}(m)}{r_{{\rm g},a}(m)}
 =e^{-2x_a(m)}.
 \label{eq:planck-ct108-lectores-exponenciales-rev3}
\end{equation}
La conjugaison échange les deux lecteurs ; l’identité de Planck est leur
équilibre unitaire.

\begin{exactbox}[title={Identité de Planck et clôture électronique}]
L'identité \(m=m_{{\rm P},a}^{\circ}\) fixe l'origine autoduale de la droite
positive des masses. La section électronique fixe une position distincte sur
cette même droite au moyen d'une généalogie propre. La
\cref{sec:ley-general-masa-accion-autonoma} construit l'opérateur universel
qui transporte échelle absolue, caractère, tours et lecture bidirectionnelle ;
la \cref{sec:cierre-electron-two-way-rev11} réalise ensuite sa réduction
centrale de rang un. Ainsi, l'unité de Planck fournit l'identité de
la coordonnée et l'électron fournit la première transition matérielle
complète sur celle-ci.
\end{exactbox}

\begin{proposition}[Réalisation généalogique du centre électronique]
\label{prop:planck-ct108-realizacion-centro-electron-rev3}
Soit
\[
 F_e^+=\mathbb R_{>0}(P_++P_-)
\]
la demi-droite positive de la fibre centrale de rang un, soit
\(\mathcal G_e\) l’espace des généalogies
électroniques enrichies et soit \(\mathcal L_M^+\) la droite positive des
masses. Un morphisme de réalisation
\[
 \mathfrak R_e:
 F_e^+\times\mathcal G_e\longrightarrow\mathcal L_M^+
\]
peut envoyer l’identité \(I=P_++P_-\) sur une masse différente de
\(m_{{\rm P},a}^{\circ}\) sans modifier le
\cref{thm:planck-ct108-recta-partida-masa-rev3} : l’identité de Planck fixe
l’origine autoduale de la carte, tandis que la généalogie électronique fixe un
déplacement sur cette carte.
\end{proposition}

\begin{proof}
Si \(\chi_e:\mathcal G_e\to\mathbb R_{>0}\) est le caractère positif construit
par le registre électronique, la règle
\[
 \mathfrak R_e(\lambda I,g)
 =\lambda\,m_{{\rm P},a}^{\circ}\chi_e(g)
\]
est positive et multiplicative dans la coordonnée généalogique. Pour la généalogie
neutre, \(\chi_e=1\), on retrouve l'origine de Planck ; pour une généalogie avec
\(\chi_e(g)\ne1\), la même identité interne occupe une position distincte sur
la droite. La clôture électronique concrète conserve sa source dans la
\cref{sec:cierre-electron-two-way-rev11}.
\end{proof}

\paragraph{Domaine de l'identité autoduale de Planck : périodicité temporelle d'ordre \(108\), réalisation circulaire déclarée et séparation de la clôture électronique}
La droite d'action pré-retour/post-retour et la périodicité temporelle
d'ordre (108) précèdent la réalisation circulaire, qui est prise comme primitive
géométrique déclarée. L'identité entre la périodicité
circulaire et les lectures de Compton, le
sélecteur \(\theta_{108}^{\rm circ}\), l'unité interne de Planck et
l'involution de la droite positive sont des résultats exacts une fois cette
réalisation fixée. La sélection gravitationnelle générale et la clôture électronique
sont construites séparément dans les sections indiquées.

\endgroup
\begingroup
\let\chapter\section
\let\section\subsection
\let\subsection\subsubsection
\let\subsubsection\paragraph
\let\clearpage\relax
\section{La loi générale avant ses réalisations par espèce}
\label{sec:ley-general-masa-accion-autonoma}
\index{masse!loi opératorielle générale}
\index{action!échelle absolue de masse}
\index{tours spectrales!loi globale de masse}

La mécanique dimensionnelle reçoit de HMT deux branches déjà construites. La première
fournit une droite d'action avec deux coordonnées orientées ; la seconde
fournit la généalogie de chaque section persistante. La masse apparaît lorsqu'on
fait agir le caractère généalogique sur une section de la droite de masse.
Cette confluence est antérieure à tout cas particulier et, par conséquent,
l'échelle \(10^{-34}\) appartient à la loi générale et non à une espèce isolée.

\subsection{Branche absolue : du déterminant local à la droite de masse}

Le lecteur déterminantiel local fixe
\begin{equation}
 \Sigma_H=\varphi_{\rm HMT}\alpha_{\rm HMT}^{16},
 \qquad \varepsilon_H:=\left\lfloor\log_{10}\Sigma_H\right\rfloor=-34 .
 \label{eq:ley-masa-sigma-menos34-autonoma}
\end{equation}
La puissance seize provient du conoyau
\(
 \mathbb Z/2\oplus\mathbb Z/2\oplus\mathbb Z/4
\)
du bloc déterminantiel. Sur la droite interne d'action
\(\mathcal L_S\), on obtient les deux sections
\begin{equation}
 \boldsymbol\hbar_{\rm pre}
 =H_5\,10^{-34}\mathcal U_S,
 \qquad
 \boldsymbol\hbar_{\rm ret}
 =\eta_{\rm ret}\,10^{-34}\mathcal U_S,
 \label{eq:ley-masa-hbar-pre-ret-autonoma}
\end{equation}
et les cartes additive et multiplicative de son dédoublement sont
\begin{equation}
 \delta_{2w}=H_5-\eta_{\rm ret},
 \qquad
 R_{\rm act}
 =\frac{\boldsymbol\hbar_{\rm pre}}
        {\boldsymbol\hbar_{\rm ret}}
 =\frac{H_5}{\eta_{\rm ret}}>0 .
 \label{eq:ley-masa-razon-accion-autonoma}
\end{equation}

Soit \(t_0\in\mathcal L_T^+\) la durée élémentaire du calendrier TPK et
\(c_{\rm int}\in\mathcal L_C^+\) sa vitesse interne. La périodicité de
108 itérations de l'endomorphisme temporel CT108 fixe la fréquence angulaire
\(\omega_0=2\pi/(108t_0)\). Les sections absolues antérieure et postérieure au
retour sont donc
\begin{equation}
 \begin{aligned}
 \mathbf m_\star^{\rm pre}
 &=\boldsymbol\hbar_{\rm pre}\,\omega_0c_{\rm int}^{-2},\\
 \mathbf m_\star^{\rm ret}
 &=\boldsymbol\hbar_{\rm ret}\,\omega_0c_{\rm int}^{-2},
 \end{aligned}
 \qquad
 \frac{\mathbf m_\star^{\rm pre}}
      {\mathbf m_\star^{\rm ret}}=R_{\rm act}.
 \label{eq:secciones-absolutas-masa-autonoma}
\end{equation}
Toutes deux appartiennent à
\(
 \mathcal L_M=\mathcal L_S\otimes\mathcal L_T^{-1}
 \otimes\mathcal L_C^{-2}
\).
Ainsi sont séparés avec précision la décade, qui fixe l'échelle absolue, et le
quotient du dédoublement bidirectionnel pré-retour/post-retour, qui
transporte la mémoire de l'évolution close.

Dans les coordonnées internes positives
\(t_0=\widehat t_0\mathcal U_T\) et
\(c_{\rm int}=\widehat c\mathcal U_C\), la réduction scalaire de toute
section dont la réalisation physique est déjà fixée prend la forme
\begin{equation}
 \mathbf m_X^\beta
 =10^{-34}\eta_{\rm ret}\,
  \frac{2\pi}{108\widehat t_0}\,\widehat c^{-2}
  a_X^{(0)}R_{\rm act}^{\kappa_X}\,\mathcal U_M .
 \label{eq:ley-escalar-general-masas-autonoma}
\end{equation}
La coordonnée \(a_X^{(0)}\) provient du caractère et de la hiérarchie harmonique ;
\(\kappa_X\) provient du lecteur de chemin. En particulier, la décade
\(-34\) appartient à la section absolue et n'est pas incorporée dans
\(\kappa_X\).

\subsection{Branche généalogique : registre, caractère et hiérarchie harmonique}

Chaque section survivante parcourt la chaîne
\begin{equation}
 \Omega_{\rm HMT}^{\rm surv}
 \longrightarrow\Gamma_{108}^{\rm enr}
 \longrightarrow\mathcal L_{108}
 \longrightarrow\sigma\in\mathbb Z^5
 \longrightarrow(\sigma,\eta)
 \longrightarrow\chi_{\rm full}(\sigma,\eta)\in\mathbb R_{>0}.
 \label{eq:cadena-genealogica-ley-masa-autonoma}
\end{equation}
Le registre conserve l'ordre des événements ; la signature conserve ses cinq
invariants entiers ; le caractère les évalue sans effacer leur provenance. Son
prolongement spectral vit dans le semi-groupe complet
\begin{equation}
 \mathfrak S=\langle90,120\rangle
 =\{90,120\}\cup\{30r:r\ge6\},
 \qquad
 s_h=q_-^h-q_+^h,
 \quad c_h=q_-^h+q_+^h,
 \label{eq:semigrupo-global-masas-autonoma}
\end{equation}
et prend la forme
\begin{equation}
 \chi_{\rm full}(\sigma,\eta)
 =\chi_0(\sigma)
  \exp\!\left(\sum_{h\in\mathfrak S}\eta_hs_h\right).
 \label{eq:caracter-completo-masas-autonoma}
\end{equation}
Les hauteurs \(90,120,180,210,240,270,360,420\) sont huit témoins
généalogiques particulièrement visibles. La loi appartient à tout
\(\mathfrak S\) ; la récurrence du second ordre de \((s_h,c_h)\) prolonge
les témoins à la hiérarchie harmonique complète.

\subsection{Opérateur universel sur les \(13\) multisections du domaine matériel}

L'atlas fournit \(81\) positions, \(56\) familles de chemins, \(13\)
multisections et \(324\) chemins orientés. Pour une multisection \(F\), soit
\(\mathcal H_F=\bigoplus_{\Gamma\in F}\mathbb C e_\Gamma\). Le caractère
basal et les tours définissent l'opérateur positif
\begin{equation}
 \widehat{\mathcal A}^{(0)}_F e_\Gamma
 =\chi_{\rm full}(\sigma_\Gamma,\eta_\Gamma)e_\Gamma.
 \label{eq:operador-caracter-basal-masas-autonoma}
\end{equation}
Les deux histoires conservées par TPK induisent des lecteurs
\(r\in\{D,\Omega\}\) : \(K_D\) agit sur le registre chronologique de
108 itérations de l'endomorphisme temporel CT108 et \(K_\Omega\), sur le mot
dodécaphasique signé. Si
\(\widehat K_F^r e_\Gamma=\kappa_r(\Gamma)e_\Gamma\), soit
\(F_s\subseteq F\) l'ensemble des chemins de la section ou du bloc physique \(s\), et
définissons
\begin{equation}
 \mathcal H_{F,s}:=\bigoplus_{\Gamma\in F_s}\mathbb C e_\Gamma,
 \qquad
 \widehat{\mathcal A}^{(0)}_{F,s}
 :=\left.\widehat{\mathcal A}^{(0)}_F\right|_{\mathcal H_{F,s}},
 \qquad
 \widehat B_{F,s}^r
 :=\left.\widehat K_F^r\right|_{\mathcal H_{F,s}}.
 \label{eq:restricciones-bloque-ley-masas-autonoma}
\end{equation}
La base diagonale des chemins utilisée par les certificats en vigueur rend invariant
\(\mathcal H_{F,s}\) sous les deux opérateurs et transforme
\(\widehat B_{F,s}^r\) en une restriction autoadjointe. Avec ces types, la loi
générale est
\begin{equation}
 \boxed{
 \widehat{\mathbf M}^{\beta,r}_{F,s}
 =\mathbf m_\star^{\rm ret}\otimes
  R_{\rm act}^{\widehat B_{F,s}^r/2}
  \widehat{\mathcal A}^{(0)}_{F,s}
  R_{\rm act}^{\widehat B_{F,s}^r/2},
 \qquad r\in\{D,\Omega\}.}
 \label{eq:ley-operatorial-general-masas-autonoma}
\end{equation}
L'écriture symétrique préserve la positivité sans présupposer la commutation. Dans la
base simultanément diagonale utilisée par les
certificats en vigueur, elle se réduit à
\begin{equation}
 \widehat{\mathbf M}^{\beta,r}_F
 =\mathbf m_\star^{\rm ret}\otimes
  \widehat{\mathcal A}^{(0)}_F
  \exp\!\left((\log R_{\rm act})\widehat K_F^r\right).
 \label{eq:ley-operatorial-diagonal-masas-autonoma}
\end{equation}
L'équation conserve, en un seul objet, l'échelle absolue, le caractère
global, toute la hiérarchie des tours et le correcteur bidirectionnel dépendant
du chemin. La carte physique de chaque secteur agit ensuite sur le spectre de
\eqref{eq:ley-operatorial-general-masas-autonoma}. Le secteur nul possède son
propre morphisme vers zéro ; les registres composés sont concaténés avant
l'évaluation de la signature ; les secteurs topologiques publient d'abord leurs invariants
de fusion et de tressage.

\begin{theorem}[Covariance et annulation de l'échelle commune]
\label{thm:covariancia-ley-general-masas-autonoma}
La loi~\eqref{eq:ley-operatorial-general-masas-autonoma} est positive et
covariante sous changement unitaire de base dans chaque fibre. Pour deux sections propres
\(X,Y\) réalisées dans la même carte de masse,
\begin{equation}
 \frac{m_Y}{m_X}
 =\frac{\chi_{\rm full}(\sigma_Y,\eta_Y)}
        {\chi_{\rm full}(\sigma_X,\eta_X)}
  R_{\rm act}^{\kappa_r(\Gamma_Y)-\kappa_r(\Gamma_X)}.
 \label{eq:cociente-ley-general-masas-autonoma}
\end{equation}
Les facteurs communs \(10^{-34}\mathcal U_S\), \(2\pi/(108t_0)\) et
\(c_{\rm int}^{-2}\) fixent la section absolue et s'annulent exactement dans
le quotient.
\end{theorem}

\begin{proof}
L'opérateur \(\widehat{\mathcal A}^{(0)}_F\) est positif parce que ses
valeurs propres appartiennent à \(\mathbb R_{>0}\) ; l'exponentielle d'un opérateur
autoadjoint est positive. La conjugaison par une matrice unitaire préserve le
spectre et, par conséquent, la loi. En divisant deux valeurs propres réalisées dans la
même carte, la section commune \(\mathbf m_\star^{\rm ret}\) se simplifie et
il reste~\eqref{eq:cociente-ley-general-masas-autonoma}.
\end{proof}

La fibre électronique centrale est de rang un et constitue le premier
corollaire scalaire de cette loi. Les autres multisections conservent le couple
d'opérateurs, leurs spectres et le morphisme physique qui sélectionne saveur,
couleur, orientation ou composition. Cette hiérarchie permet de publier ensemble le
résultat algébrique, sa réalisation dimensionnelle et la confrontation externe sans
confondre leurs domaines.

\endgroup
\begingroup
\let\chapter\section
\let\section\subsection
\let\subsection\subsubsection
\let\subsubsection\paragraph
\let\clearpage\relax
% Vista científica exacta materializada desde una rama condicional.
% Fuente: source/demostraciones_primarias/09_06h_atlas_familias_rev6.tex; rama: HMTIncludeAtlasStructural.
% No modifica el contenido matemático de la rama de origen.
\chapter[Classification nonadique des familles]{Classification nonadique des familles de trajectoires}
\label{chap:atlas-genealogia-masas}
\index{atlas nonadique}
\index{famille de chemin}
\index{multisection}
\index{secteur nul}

L'atlas ne commence pas par une liste de particules ni par une table de masses.
Il commence par un support fini orienté et par les lecteurs qui conservent
l'histoire de ses chemins. Sa fonction est de répondre à une question antérieure à
toute comparaison métrologique : quelles classes de persistance la géométrie nonadique peut-elle
soutenir et quelle information doit être conservée pour les distinguer ?

La réponse comporte trois cardinaux, mais non trois inventaires rivaux :
\[
 81\ \text{cellules},\qquad
 56\ \text{familles de routes},\qquad
 13\ \text{multisections famille--niveau}.
\]
Le premier nombre compte des positions orientées ; le deuxième identifie les
positions qui publient la même signature finie de chemin ; le troisième réunit les
positions qui partagent famille et profondeur par rapport au centre. Aucun
d'eux ne compte des masses connues, des noms d'espèces ou des lignes d'une
comparaison externe.

\section{La géométrie complète des \(81\) cellules}

Soit
\[
 I_9=\{1,\ldots,9\},
 \qquad
 \mathcal C_{81}=I_9\times I_9.
\]
Les coordonnées \((r,c)\) conservent les deux directions d'APP\@. La parité
de chaque coordonnée définit quatre secteurs internes :
\[
 \operatorname{fam}(r,c)=
 \begin{cases}
  \ell^- ,& r\equiv1,\ c\equiv1\pmod2,\\
  \nu ,& r\equiv1,\ c\equiv0\pmod2,\\
  d ,& r\equiv0,\ c\equiv1\pmod2,\\
  u ,& r\equiv0,\ c\equiv0\pmod2.
 \end{cases}
 \tag{A1}\label{eq:familias-paridad-atlas}
\]
La profondeur dans l'atlas n'est pas ajoutée comme étiquette indépendante. C'est la
distance de Tchebychev au centre :
\[
 G(r,c)=\max\{|r-5|,|c-5|\}.
 \tag{A2}\label{eq:generacion-chebyshev-atlas}
\]
Par conséquent, famille et niveau proviennent de la position elle-même. En particulier,
le centre n'est pas obtenu en sélectionnant une espèce de l'extérieur : c'est l'unique
cellule de profondeur zéro.

\begin{theorem}[Recensement des familles et des générations de l'atlas]
\label{thm:censo-atlas-81}
Les applications \eqref{eq:familias-paridad-atlas} et
\eqref{eq:generacion-chebyshev-atlas} décomposent l'atlas de manière
disjointe comme
\[
 81=45_{\rm leptones}+36_{\rm quarks}
   =25_{\ell^-}+20_{\nu}+20_{d}+16_{u}.
 \tag{A3}\label{eq:censo-81-cuatro-familias}
\]
En raffinant par \(G\), on obtient exactement les treize multisections
\[
\begin{array}{c|c|c}
\text{famille}&\text{niveaux}&\text{cardinaux}\\ \hline
\ell^-&G0,G2,G4&1,8,16\\
\nu&G1,G2,G3,G4&2,4,6,8\\
d&G1,G2,G3,G4&2,4,6,8\\
u&G1,G3&4,12.
\end{array}
\tag{A4}\label{eq:trece-multisecciones-atlas}
\]
La cellule \((5,5)\) est l'unique composante de
\(\mathfrak S_{\ell^-,0}\).
\end{theorem}

\begin{proof}
Dans \(I_9\), il y a cinq coordonnées impaires et quatre paires. Les quatre classes
de parité ont donc pour cardinaux
\[
 5\cdot5=25,\qquad
 5\cdot4=20,\qquad
 4\cdot5=20,\qquad
 4\cdot4=16,
\]
ce qui prouve \eqref{eq:censo-81-cuatro-familias} ; les deux premières
classes réunissent \(45\) cellules et les deux dernières, \(36\).

Pour le raffinement, nous faisons croître des carrés centrés en \((5,5)\).
Dans la classe impair--impair, il y a une cellule au rayon zéro, \(3^2-1=8\)
dans la couche de rayon deux et \(5^2-3^2=16\) dans celle de rayon quatre.
Dans la classe impair--pair, les rectangles accumulés jusqu'aux rayons
\(1,2,3,4\) contiennent, respectivement,
\[
 1\cdot2=2,\qquad 3\cdot2=6,\qquad
 3\cdot4=12,\qquad5\cdot4=20
\]
cellules ; leurs différences successives sont \(2,4,6,8\). La classe
pair--impair donne le même calcul avec les coordonnées échangées. Dans la
classe pair--pair apparaissent \(2^2=4\) cellules au rayon un et
\(4^2-2^2=12\) au rayon trois. Ces couches sont disjointes et épuisent
\(I_9^2\), de sorte qu'il ne reste aucune autre multisection.
\end{proof}

Le mot \emph{niveau} désigne dans ce chapitre la profondeur combinatoire
\(G\). La lecture physique ultérieure conserve cette géométrie, mais ne la
définit pas. On évite ainsi de transformer les noms des particules en prémisses
du recensement qui doit les expliquer.

\section{Inversion centrale, \(41\) orbites et unicité du centre}

L'inversion de l'atlas est
\[
 \rho_{\rm c}(r,c)=(10-r,10-c).
 \tag{A5}\label{eq:inversion-central-atlas-rev6}
\]
Comme \(10-r\) a la même parité que \(r\), l'inversion conserve les
quatre familles. En outre,
\[
 |(10-r)-5|=|r-5|,
 \qquad
 |(10-c)-5|=|c-5|,
\]
et, par conséquent, conserve le niveau \(G\). Chacune des treize
multisections est \(\rho_{\rm c}\)-stable.

\begin{proposition}[Orbites et exception centrale]
\label{prop:orbitas-centro-atlas-rev6}
L'action de \(\rho_{\rm c}\) sur \(\mathcal C_{81}\) possède exactement
\(41\) orbites. Son unique point fixe est \((5,5)\). Par conséquent, la
multisection \(\mathfrak S_{\ell^-,0}\) admet l'unique sélecteur ponctuel
\(\rho_{\rm c}\)-équivariant ; dans les douze autres, l'objet naturel est
l'orbite ou la multisection complète.
\end{proposition}

\begin{proof}
L'équation
\[
 (r,c)=\rho_{\rm c}(r,c)
\]
équivaut à \(2r=2c=10\), donc \(r=c=5\). Les quatre-vingts autres cellules se
regroupent en quarante couples, d'où \(1+40=41\) orbites.

Soit \(y\) une multisection non centrale et supposons qu'un sélecteur ponctuel
équivariant choisisse \(a(y)\in\mathfrak S_y\). Comme \(\rho_{\rm c}\) agit
trivialement sur le type famille--niveau, l'équivariance imposerait
\[
 a(y)=a(\rho_{\rm c}y)=\rho_{\rm c}a(y).
\]
Le point choisi devrait être fixe, mais l'unique point fixe est dans
\(\mathfrak S_{\ell^-,0}\). La contradiction démontre l'affirmation.
\end{proof}

Cette proposition explique pourquoi une espèce interne n'est pas, en général, une
case isolée. L'électron central est exceptionnel parce que sa multisection
se réduit à un point. Dans les autres familles, l'orientation conjuguée fait
partie de l'identité qui doit être conservée.

\section{Les \(56\) familles de chemins}

Chaque cellule porte l'ombre finie de son histoire de transport :
\[
 \mathcal R(r,c)=
 \bigl(
 n_A,\ s_C,\ k_{\Delta_4},\
 \nu_{120},\ \nu_{270},\ \tau_{12}
 \bigr)_{r,c}.
 \tag{A6}\label{eq:lector-familia-ruta-atlas}
\]
Les cinq premières coordonnées sont entières ; \(\tau_{12}\) est le mot
dodécaphasique orienté. Deux cellules appartiennent à la même famille de chemins
si et seulement si elles ont le même sextuplet \(\mathcal R\).

\begin{theorem}[Recensement des familles de chemins et mémoire minimale]
\label{thm:censo-56-familias-ruta-rev6}
L'image de \(\mathcal R\) contient exactement \(56\) éléments.
L'histogramme des tailles de ses fibres est
\[
 41\times1,\qquad
 10\times2,\qquad
 2\times3,\qquad
 1\times4,\qquad
 2\times5.
 \tag{A7}\label{eq:histograma-fibras-ruta}
\]
En particulier,
\[
 41+10+2+1+2=56,
\]
\[
 41\cdot1+10\cdot2+2\cdot3+1\cdot4+2\cdot5=81.
 \tag{A8}\label{eq:doble-censo-rutas}
\]
L'ordre lexicographique étant fixé dans chaque fibre, le rang
\[
 \kappa(r,c)\in\{0,1,2,3,4\}
\]
rend injective l'application
\[
 (r,c)\longmapsto\bigl(\mathcal R(r,c),\kappa(r,c)\bigr).
 \tag{A9}\label{eq:lift-memoria-cinco-atlas}
\]
Cinq est la taille minimale d'un alphabet uniforme capable d'effectuer ce
relèvement.
\end{theorem}

\begin{proof}
L'énumération exhaustive de \(I_9^2\) regroupe les quatre-vingt-un sextuplets
\eqref{eq:lector-familia-ruta-atlas} par égalité coordonnée par coordonnée.
Elle produit les cinq tailles de fibre de
\eqref{eq:histograma-fibras-ruta} ; les deux identités
\eqref{eq:doble-censo-rutas} vérifient simultanément le nombre de classes
et la couverture des quatre-vingt-une cellules.

Le rang lexicographique distingue les éléments de chaque fibre, de sorte que
\eqref{eq:lift-memoria-cinco-atlas} est injective. Comme il existe deux fibres
de cardinal cinq, tout alphabet commun comportant moins de cinq symboles
identifierait deux cellules de l'une d'elles par le principe des tiroirs.
\end{proof}

Le rang \(\kappa\) est une mémoire de position au sein d'une fibre ; il n'est
ni charge, ni masse, ni couleur. Sa fonction est de retrouver la cellule que la signature visible
du chemin avait identifiée à d'autres. Le nombre \(56\) compte les classes du
sextuplet \(\mathcal R\), tandis que le nombre \(13\) compte les classes du couple
\((\operatorname{fam},G)\). Ce sont des quotients différents du même atlas.

\paragraph{Quotients cellulaire et orbital de l'atlas : domaines et applications de projection.}
La projection brute CT--108 et la classification scientifique lisent des données
distinctes des mêmes quatre-vingt-une cellules. Si
\(q_{\rm CT}\) conserve seulement la forme du mot tritique
dodécaphasique, tandis que \(\mathcal R\) conserve le sextuplet de chemin et
\(q=(\operatorname{fam},G)\) conserve famille et profondeur, le tableau
typé est
\[
\begin{array}{rcll}
 \mathcal C_{81}
 &\xrightarrow{\ q_{\rm CT}\ }&
 \mathcal W^{\rm raw}_{7}
 & \text{formes brutes CT--108},\\[2mm]
 \mathcal C_{81}
 &\xrightarrow{\ \mathcal R\ }&
 \Sigma^{\rm ruta}_{56}
 & \text{familles de routes},\\[2mm]
 \mathcal C_{81}
 &\xrightarrow{\ q\ }&
 \Sigma_{13}
 & \text{multisections famille--niveau}.
\end{array}
\tag{A9a}\label{eq:dos-cocientes-atlas-rev8}
\]
Les sept fibres brutes ont pour multiplicités
\[
 54,2,2,2,7,7,7,
 \qquad54+3\cdot2+3\cdot7=81.
\]
Par conséquent, le recensement archéologique \(81\to7\) ne remplace pas l'atlas
\(81/56/13\). L'écriture abrégée \(81/56/13\) n'affirme pas non plus qu'il existe
une flèche canonique \(56\to13\) : certaines fibres de \(\mathcal R\) traversent
plus d'un couple famille--niveau. Les applications \(\mathcal R\) et \(q\) sont des
quotients parallèles et conservent des invariants différents. Dans le secteur quark,
la convention active reste
\[
 \operatorname{col}(r,c)=r-c\pmod3;
\]
l'écriture historique \(c-r\) est sa convention conjuguée, non une autre couleur ni
un autre quotient de l'atlas.

\section{Distinction entre cellule, famille de chemins et multisection}

Définissons
\[
 q:\mathcal C_{81}\longrightarrow\Sigma_{13},
 \qquad
 q(r,c)=\bigl(\operatorname{fam}(r,c),G(r,c)\bigr),
 \tag{A10}\label{eq:cociente-trece-multisecciones}
\]
et, pour \(y\in\Sigma_{13}\),
\[
 \mathfrak S_y^{\rm atlas}
 =q^{-1}(y)
 =\{(r,c)\in\mathcal C_{81}:q(r,c)=y\}.
 \tag{A11}\label{eq:multiseccion-atlas-rev6}
\]
Les quatre niveaux de lecture sont alors séparés :

\begin{enumerate}[label=(\roman*)]
 \item une \emph{cellule} est une position orientée \((r,c)\), avec toutes ses
 données locales ;
 \item une \emph{famille de chemins} est une fibre de \(\mathcal R\) : elle conserve
 une même signature finie de transport ;
 \item une \emph{multisection} est une fibre de \(q\) : elle conserve famille,
 profondeur et totalité de ses orientations conjuguées ;
 \item une \emph{représentation physique} est la réalisation ultérieure d'une
 multisection enrichie dans un espace d'états ; ce n'est pas une autre partition du
 tableau ni le choix arbitraire de l'une de ses cellules.
\end{enumerate}

Cette distinction donne un contenu mathématique au mot \emph{espèce}. Le type
interne est la multisection. Une particule concrète est une section
persistante réalisée en elle. La représentation physique conserve l'action,
l'orientation et les observables du type ; elle ne revient pas en arrière pour redéfinir
l'atlas à partir d'une masse déjà connue.

\section{Couleur résiduelle comme orientation ternaire}

Dans le secteur quark, formé des \(36\) cellules avec \(r\) pair, nous définissons
\[
 \operatorname{col}(r,c)=r-c\pmod3,
 \tag{A12}\label{eq:color-residual-atlas-rev6}
\]
avec le dictionnaire
\[
 0\longmapsto R,\qquad1\longmapsto G,\qquad2\longmapsto B.
\]

\begin{proposition}[Bilan de couleur résiduelle]
\label{prop:color-atlas}
Les trente-six cellules de quark se répartissent exactement comme
\[
 36=12_R+12_G+12_B.
 \tag{A13}\label{eq:balance-color-atlas-rev6}
\]
L'inversion centrale fixe \(R\) et échange \(G\) et \(B\).
\end{proposition}

\begin{proof}
L'une quelconque des quatre valeurs paires de \(r\) étant fixée, lorsqu'on parcourt
\(c=1,\ldots,9\), chaque classe modulo trois apparaît trois fois. Chaque résidu
a donc \(4\cdot3=12\) représentants. En outre,
\[
 \operatorname{col}\bigl(\rho_{\rm c}(r,c)\bigr)
 =(10-r)-(10-c)
 =-(r-c)\pmod3.
\]
La négation modulaire fixe la classe zéro et échange les classes un et deux.
\end{proof}

Le tritique \(R,G,B\) est donc un observable exact de l'atlas et non une
décoration ajoutée à une table de quarks. La représentation de couleur se
construit ensuite sur ce support ternaire ; le recensement \(12+12+12\) est sa
géométrie préalable.

\section{De l'extension survivante à la signature}

L'atlas classe des coupes finies. L'identité d'une particule requiert
l'histoire profonde dont cette coupe provient. Soit
\[
 \Omega_{\rm HMT}^{\rm surv}
 =\varprojlim_m\mathcal S_m
\]
le système inverse des états survivants. La projection sur l'atlas \(\pi_{81}\) et l'ombre
du chemin forment la chaîne
\[
 \Omega_{\rm HMT}^{\rm surv}
 \xrightarrow{\ \operatorname{sh}_{108}\ }
 \Gamma_{108}^{\rm enr}
 \xrightarrow{\ \mathcal E\ }
 \mathcal L_{108}
 \xrightarrow{\ L_{\rm tick}\ }
 \mathbb Z^5.
 \tag{A14}\label{eq:cadena-extension-firma-atlas-rev6}
\]
Ici, \(\operatorname{sh}_{108}\) publie un chemin de \(108\) pas sans
oublier feuillet, orientation, frontière ni prédécesseurs ; \(\mathcal E\) forme le
registre dodécaphasique ; et \(L_{\rm tick}\) extrait
\[
 v(x)=
 \bigl(n_A,s_C,k_{\Delta_4},\nu_{120},\nu_{270}\bigr).
 \tag{A15}\label{eq:firma-cinco-extension-rev6}
\]
Aucune de ces trois flèches ne consulte une masse, une unité ou le nom
d'une particule.

La raison de conserver le registre avant la projection se voit dans ses douze
événements orientés \(e_0,\ldots,e_{11}\). Pour \(0\leq i\leq5\), soit
\[
 p_i=e_i+e_{i+6},
 \qquad
 a_i=e_i-e_{i+6},
 \qquad
 s_C=\sum_{i=0}^{5}p_i,
\]
et, pour \(0\leq i<5\),
\[
 M_i=p_i-p_5.
\]
Alors
\[
 p_5=\frac{s_C-\sum_{i=0}^{4}M_i}{6},
 \qquad
 p_i=p_5+M_i,
\]
\[
 e_i=\frac{p_i+a_i}{2},
 \qquad
 e_{i+6}=\frac{p_i-a_i}{2}.
 \tag{A16}\label{eq:reconstruccion-registro-doce}
\]
L’identité
\[
 1+5+6=12
\]
est ainsi une reconstruction réversible : un total, cinq différences et six
orientations retrouvent les douze événements. Projeter sur la signature de cinq
entiers avant de composer les histoires éliminerait précisément le placement et
l'orientation qui permettent de décider si deux chemins peuvent être composés.

Le quotient \(q\) se relève à l'histoire profonde au moyen de la projection sur l'atlas
\(\pi_{81}:\Omega_{\rm HMT}^{\rm surv}\to\mathcal C_{81}\) :
\[
 \widetilde{\mathfrak S}_y
 =\{x\in\Omega_{\rm HMT}^{\rm surv}:q(\pi_{81}x)=y\}.
 \tag{A17}\label{eq:multiseccion-enriquecida-rev6}
\]
C'est la multisection enrichie. Son ombre appartient à l'atlas ; son chemin
appartient à \(\Gamma_{108}^{\rm enr}\) ; sa mémoire complète appartient au
registre ; et son évaluation n'est appliquée qu'à la fin. La généalogie n'est pas une
note ajoutée à la valeur : c'est l'objet qui rend la valeur possible et conserve les
autres lectures de la même section.

\section{Bifurcation du secteur nul avant le caractère positif}

Le caractère de masse agit sur la signature entière et a un codomaine positif :
\[
 \chi_{\rm mass}:\mathbb Z^5\longrightarrow\mathbb R_{>0}.
 \tag{A18}\label{eq:caracter-positivo-atlas-rev6}
\]
Par définition, il n'existe pas de signature \(v\) avec
\(\chi_{\rm mass}(v)=0\). Le secteur nul n'est donc pas une
quatorzième multisection ni la limite dégénérée des treize précédentes.

Son support algébrique appartient à la branche scindée
\[
 \mathcal A_{\rm sp}
 =\mathbb R[\mathsf R]/(\mathsf R^2-1),
 \qquad
 P_\pm=\frac{1\pm\mathsf R}{2}.
\]
Tout élément s'écrit de manière unique comme
\[
 z=r_+P_++r_-P_-,
 \qquad
 N_{\rm sp}(z)=r_+r_-,
\]
et les deux droites
\[
 \mathbb R P_+,\qquad\mathbb R P_-
\]
sont nulles parce que \(N_{\rm sp}(\lambda P_\pm)=0\). Une section nulle est un
chemin ouvert transporté dans l'une de ces directions. L'intérieur matériel
requiert le couplage des deux, pour lequel \(r_+r_->0\).

Ainsi sont typées deux branches :
\[
 \text{persistance matérielle}
 \longrightarrow\mathbb Z^5
 \xrightarrow{\chi_{\rm mass}}\mathbb R_{>0},
\]
\[
 \text{propagation nulle}
 \longrightarrow\partial_{\rm null}\mathcal C_{\rm sp}^{+}.
\]
Les réalisations photonique et de jauge lisent ensuite la seconde branche. Sa
séparation du caractère positif empêche d'effacer photon et gluon parce qu'ils n'apparaissent pas
dans une liste de masses, et évite en même temps de fabriquer une masse nulle au moyen d'un
caractère dont le codomaine l'exclut.

\section{Système fini de classification et mémoire des trajectoires}

Les trois quotients du chapitre remplissent des fonctions complémentaires :
\[
 \mathcal C_{81}
 \xrightarrow{\ \mathcal R\ }\mathcal R_{56},
 \qquad
 \mathcal C_{81}
 \xrightarrow{\ q\ }\Sigma_{13},
 \qquad
 \mathcal C_{81}
 \xrightarrow{\ /\langle\rho_{\rm c}\rangle\ }
 \mathcal O_{41}.
\]
Le premier conserve la signature du chemin ; le deuxième conserve famille et niveau ;
le troisième conserve la conjugaison centrale. Avec la couleur résiduelle et le
relèvement \(\kappa\), ils forment un système fini capable de reconstruire
quelle information est visible, laquelle est identifiée et laquelle doit subsister comme
mémoire.

La chaîne causale de l'atlas est donc fixée :
\[
 \boxed{
 \begin{gathered}
 \text{extension survivante}
 \longrightarrow\text{route enrichie}
 \longrightarrow\text{registre}
 \longrightarrow\text{signature}\\
 \longrightarrow\text{caractère}
 \longrightarrow\text{lecture dimensionnelle}
 \end{gathered}}
 \tag{A19}\label{eq:cadena-rectora-atlas-rev6}
\]
Les quatre-vingt-une cellules constituent le domaine local ; les cinquante-six
familles expriment des coïncidences de chemins ; les treize multisections sont les
types internes ; les quarante-et-une orbites conservent la conjugaison ; et le
centre \((5,5)\) est l'unique exception ponctuelle. La masse apparaît ensuite
comme lecture d'une persistance déjà construite. L'atlas n'est pas une table
abrégée de résultats : c'est la structure qui explique pourquoi une table peut
exister.

\endgroup
\begingroup
\let\chapter\section
\let\section\subsection
\let\subsection\subsubsection
\let\subsubsection\paragraph
\let\clearpage\relax
\section{Des germes APP au registre de mélange et de masse}
\label{sec:semillas-registro-mezcla-autonoma}
\index{germe!factorisation somme--produit}
\index{mélange!provenance dans les germes}

La loi physique reçoit un domaine qui a déjà conservé l'incompatibilité des
deux évaluations APP. Si
\(\mathcal S_+\) et \(\mathcal S_\times\) désignent les germes orientés des
feuillets additif et multiplicatif, alors
\begin{equation}
 \Omega_{\rm seed}=\mathcal S_+\times\mathcal S_\times,
 \qquad
 |\Omega_{\rm seed}|=324^2=104,976.
 \label{eq:dominio-semillas-md-autonoma}
\end{equation}
La factorisation exacte de l'émetteur produit
\begin{equation}
 104,976
 \longrightarrow18\times26=468\ U_6
 \longrightarrow243\ w_6\subset\mathbb F_3^6,
 \qquad |\mathbb F_3^6|=729.
 \label{eq:embudo-semillas-md-autonoma}
\end{equation}
Les quatre cardinaux comptent des objets distincts : conditions initiales de
deux feuillets, émissions décimales, mots ternaires publiés et espace ambiant
algébrique. Chaque germe conserve, outre son émission, multiplicité,
diagonale, feuillet, orbite, signature d'émission, résidus modulo trois et modulo
neuf, rotations, changements de feuillet, extrémité, déplacements et comptages
cardinaux. Ces données constituent le vocabulaire du registre ultérieur.

La factorisation organise six microfibres de cardinal \(1008\), six plans
de Fano de phase, trente-six fibres de cardinal \(28\) et la bande
\(1944=27\cdot72\), diagonalisée par \((\mathbb Z/3)^3\). En moyennant le
mot nonadique du TPK, les projecteurs conditionnels des deux feuillets
définissent le noyau canonique
\begin{equation}
 K_{\rm cat}=\frac13(I+E_++E_\times),
 \qquad
 \operatorname{Spec}(K_{\rm cat})
 =\left\{1^{[1]},
 \left(\frac23\right)^{[42]},
 \left(\frac13\right)^{[425]}\right\}.
 \label{eq:kernel-canonico-semillas-md-autonoma}
\end{equation}
La phase TPK sélectionne le renouvellement actif et conserve les neuf bandes
spatiales ; la moyenne~\eqref{eq:kernel-canonico-semillas-md-autonoma}
publie l'ombre stochastique du calendrier complet.

\subsection{Bigraduation des germes et \(6\) plans d'incidence}

Dans chacune des six microfibres de cardinalité \(1008\), le graphe
biparti conserve \(54\) sommets additifs et \(56\) multiplicatifs. Les
premiers se décomposent en \(36\) sommets de degré \(24\) et \(18\) de
degré \(8\) ; les seconds ont le degré \(18\). Le décompte des incidences
coïncide pour les deux feuillets :
\begin{equation}
 1008=36\cdot24+18\cdot8=56\cdot18.
 \label{eq:bidegraduacion-microfibra-masas-autonoma}
\end{equation}
Cette égalité fixe la bigraduation de l'émetteur avant sa projection sur
les familles de chemins.

Les six mots
\begin{equation}
 001112,\quad010211,\quad100121,\quad112001,\quad121100,\quad211010
 \label{eq:seis-palabras-microfibra-masas-autonoma}
\end{equation}
forment une orbite de \(D_3\simeq S_3\) et possèdent le profil
\(432+4\cdot144=1008\). Le quotient \(1008=36\cdot28\) organise sept
blocs
\[
 Q=\{B_3,B_6,B_9,P_1,P_2,P_3,f\},
\]
où \((B_3,B_6,B_9)\) sont les tétrades transversales et
\((P_1,P_2,P_3,f)\), les quatre tétrades latérales. Pour l'un des six
mots \(a\), soit
\[
 \mathcal R_a
 :=\{(p,t)\in\mathcal S_+\times\mathcal S_\times:
       U_6(p,t)\text{ publie le mot }a\}
\]
l'ensemble de ses mille huit chemins, et posons
\(R_9:=\mathbb Z/9\mathbb Z\). La carte orientée étant fixée, la projection
de phase a pour type
\begin{equation}
 q_\beta:\mathcal R_a\longrightarrow\binom{R_9}{2},
 \qquad
 q_\beta(p,t)
 =\{x-5\beta(d_t),x+5\beta(d_t)\}\subset\mathbb Z/9\mathbb Z,
 \qquad x=5(i_p-j_p),
 \label{eq:proyeccion-fano-semillas-masas-autonoma}
\end{equation}
avec \((\beta(N),\beta(E),\beta(S),\beta(O))=(1,2,3,4)\). Son codomaine est
l'ensemble des trente-six couples non ordonnés de \(R_9\). Chaque
\(X\in Q\) est une tétrade d'orbites ;
son support de chemins et sa coupe sur une fibre sont définis par
\[
 \operatorname{Supp}_{\rm rutas}(X)
 :=\bigcup_{\mathcal O\in X}\mathcal O,
 \qquad
 X_b:=\operatorname{Supp}_{\rm rutas}(X)\cap q_\beta^{-1}(b),
 \qquad b\in\binom{R_9}{2}.
\]
Dans le canal multiplicatif, l'observable source est
\begin{equation}
 \rho:\mathcal S_\times\longrightarrow R_9,
 \qquad
 \rho(i,j,d)=(i+1)(j+1)\pmod9.
 \label{eq:residuo-producto-fano-semillas-rev3}
\end{equation}
Si \(\pi_\times:\mathcal R_a\to\mathcal S_\times\),
\(\pi_\times(p,t)=t\), est la seconde projection, l'observable qui agit
sur les chemins est
\begin{equation}
 \rho_\times:=\rho\circ\pi_\times:\mathcal R_a\longrightarrow R_9,
 \qquad
 \rho_\times(p,t)=(i_t+1)(j_t+1)\pmod9.
 \label{eq:residuo-producto-rutas-fano-semillas-rev3}
\end{equation}
Dans ce qui suit, \(\rho(X)\) abrège le multiensemble des valeurs de
\(\rho_\times\) sur \(X_b\), qui est indépendant de la fibre \(b\). Ses
profils sur les tétrades sont
\[
 \rho(B_3)=3^4,
 \qquad
 \rho(B_6)=6^4,
 \qquad
 \rho(B_9)=0^4,
\]
et, pour les trois tétrades latérales,
\begin{equation}
 \begin{array}{c|c|c}
 r&\text{signature active de }P_r&R(P_r)\\
 \hline
 1&933&\{5,7\}\\
 2&393&\{4,8\}\\
 3&339&\{1,2\}.
 \end{array}
 \label{eq:perfiles-laterales-fano-semillas-rev3}
\end{equation}
Chaque résidu de \(R(P_r)\) apparaît deux fois parmi les quatre chemins de
\((P_r)_b\).

\begin{theorem}[Morphisme d'incidence des germes et six plans de Fano]
\label{thm:semillas-seis-fano-fase-rev3}
Pour \(\delta\in R_9\), \(h\in\{3,6,0\}\) ---correspondant,
respectivement, à \(B_3,B_6,B_9\)--- et \(r\in\{1,2,3\}\), soit
\begin{equation}
 N_\delta(B_h,P_r)
 :=\sum_{b\in\binom{R_9}{2}}
 \#\left\{(x,y)\in(B_h)_b\times(P_r)_b:
 \rho_\times(y)-\rho_\times(x)=\delta\right\}.
 \label{eq:conteo-incidencia-fano-semillas-rev3}
\end{equation}
Alors
\begin{equation}
 N_\delta(B_h,P_r)=
 \begin{cases}
 288,&h+\delta\in R(P_r),\\
 0,&h+\delta\notin R(P_r).
 \end{cases}
 \label{eq:conteo-288-cero-fano-semillas-rev3}
\end{equation}
Si \(\delta\in R_9^\times\), pour chaque \(B_h\) il existe un unique \(P_r\)
avec comptage \(288\) ; on obtient donc une bijection
\[
 \theta_\delta:\{B_3,B_6,B_9\}
 \xrightarrow{\sim}\{P_1,P_2,P_3\}.
\]
En écrivant l'image de \((B_3,B_6,B_9)\) au moyen de ses indices, les six
bijections sont
\begin{equation}
 \begin{array}{c|c@{\qquad}c|c}
 \delta&\theta_\delta&\delta&\theta_\delta\\
 \hline
 1&(2,1,3)&5&(2,3,1)\\
 2&(1,2,3)&7&(3,2,1)\\
 4&(1,3,2)&8&(3,1,2).
 \end{array}
 \label{eq:tabla-theta-delta-fano-semillas-rev3}
\end{equation}
Soit \(P=\{P_1,P_2,P_3\}\) et, pour \(p\in P\), soit
\[
 \mu(p):=\bigl\{\{f,p\},P\setminus\{p\}\bigr\}
\]
le couplage parfait associé sur les quatre tétrades latérales.
Alors
\begin{equation}
 \mathcal D_\delta
 :=\bigl\{\{B_3,B_6,B_9\}\bigr\}
 \cup
 \bigl\{\{B,q,q'\}:B\in\{B_3,B_6,B_9\},
       \{q,q'\}\in\mu(\theta_\delta(B))\bigr\}
 \label{eq:sistema-fano-semillas-rev3}
\end{equation}
est un système de Steiner \(\operatorname{S}(2,3,7)\). Les six unités
produisent les six plans de Fano sur ces sept blocs qui contiennent la
droite transversale \(\{B_3,B_6,B_9\}\).
\end{theorem}

\begin{proof}
Fixons une fibre \(b\). Dans \((B_h)_b\), il existe quatre choix de
résidu \(h\). Si \(h+\delta\in R(P_r)\), le profil
\eqref{eq:perfiles-laterales-fano-semillas-rev3} fournit exactement deux
choix dans \((P_r)_b\) avec ce résidu ; s'il n'appartient pas au support, il
n'en fournit aucun. Comme il y a trente-six fibres,
\(N_\delta=36\cdot4\cdot2=288\) dans le premier cas et zéro dans le second.
Les trois supports latéraux partitionnent \(R_9^\times\), de sorte qu'une unité
\(\delta\) détermine un unique latéral pour chaque bloc transversal. La lecture
directe de ces supports donne le tableau
\eqref{eq:tabla-theta-delta-fano-semillas-rev3} et épuise les \(3!=6\)
bijections.

La droite transversale couvre les trois couples entre \(B_3,B_6,B_9\). Pour un
\(B\) fixé, les deux arêtes du couplage \(\mu(\theta_\delta(B))\)
couvrent une fois les quatre points latéraux. Les trois couplages sont
distincts et forment la factorisation de \(K_4\) en ses trois couplages
parfaits ; par conséquent, chaque
couple latéral apparaît exactement une fois. Les sept triplets contiennent ainsi
chacun des vingt-et-un couples de points une fois, ce qui est la propriété
\(\operatorname{S}(2,3,7)\). Réciproquement, tout plan de Fano contenant
la droite transversale associe à chaque \(B\) un couplage parfait distinct
des deux autres et retrouve l'une des six bijections précédentes.
\end{proof}

La correspondance d'ensembles \(\delta\mapsto\theta_\delta\) n'est ni une
action ni un homomorphisme : \(R_9^\times\simeq C_6\) est abélien pour le
produit, tandis que \(S_3\) ne l'est pas, et les unités ne forment pas un sous-groupe
additif. Le déplacement \(\delta\) indexe six translations des profils
du produit.

Le même calcul et le même tableau se vérifient dans les six microfibres. Il faut
néanmoins conserver la différence entre support et rôle : le support de
\(f\) coïncide avec l'un des supports latéraux et, par conséquent, l'observable
\(\rho\) ne retrouve pas à lui seul la décomposition \(P\sqcup\{f\}\). Cette
distinction exige la signature additive déjà déclarée dans l'état de germe.

\begin{statusbox}[title={Statut et portée de l'incidence de Fano}]
Le \cref{thm:semillas-seis-fano-fase-rev3} est un théorème fini une fois
fixés \(q_\beta\), l'orientation et l'ordre des chiffres : il définit six jauges
exactes d'incidence sur les blocs. Il organise phase et incidence dans le domaine
des germes ; il ne sélectionne ni particule, ni masse, ni chemin. Sa réalisation physique
s'effectue ensuite au moyen des morphismes du registre enrichi.
\end{statusbox}

Le passage à la physique conserve la chaîne
\begin{equation}
 \begin{aligned}
 \text{germe raffiné}
 &\longrightarrow(\text{signatures }+,\times;U_6;w_6;
                   \text{mémoire de feuillet})\\
 &\longrightarrow\text{registre CT108}
 \longrightarrow(\tau_{12},D_k,\Omega,P_6,\nu_{120},\nu_{270})\\
 &\longrightarrow\text{fibre et multisection}
 \longrightarrow(B_D,B_\Omega)
 \longrightarrow\text{transport CKM/PMNS}\\
 &\longrightarrow\text{caractère et tours}
 \longrightarrow\text{section de masse}.
 \end{aligned}
 \label{eq:cadena-semilla-mezcla-masa-autonoma}
\end{equation}
CKM réalise un changement de base des opérateurs de quark,
\begin{equation}
 B_d^{(u)}=V_{\rm CKM}B_dV_{\rm CKM}^*,
 \qquad
 R_{\rm act}^{B_d^{(u)}}
 =V_{\rm CKM}R_{\rm act}^{B_d}V_{\rm CKM}^*,
 \label{eq:ckm-transporte-semillas-autonoma}
\end{equation}
et préserve spectre et déterminant. Le mélange spécifie la représentation
de l'opérateur dans une base de saveur ; le registre conserve les valeurs propres et leur
généalogie.

\subsection{Lecteur neutrino contenu dans le registre}

Le secteur neutre dispose d'un extracteur entier construit exclusivement
avec des données du registre. Pour chaque bloc dodécaphasique \(E_m\), nous définissons
\begin{equation}
 T(m)=\sum_{t\in E_m}(-1)^{\kappa(t)}\varepsilon(t)
       \mathbf1_{\{d(t)=9\}},
 \qquad
 s_C^{(m)}(\nu)=T(m)-T(m+6),
 \label{eq:lector-neutrino-ledger-autonoma}
\end{equation}
et
\begin{equation}
 \tau_m^\nu=\operatorname{sgn}s_C^{(m)}(\nu).
 \label{eq:trit-neutrino-ledger-autonoma}
\end{equation}
Les quatre triangles de saut quatre produisent
\begin{equation}
 \nu_{120}=\sum_{j=1}^{4}
 \operatorname{sgn}\!\left(\sum_{m\in\mathcal T_j}\tau_m^\nu\right),
 \qquad
 \nu_{270}=\sum_{m=1}^{12}\tau_m^\nu\tau_{m+3}^\nu.
 \label{eq:invariantes-neutrino-ledger-autonoma}
\end{equation}
L'inversion bidirectionnelle transforme
\(\nu_{120}\mapsto-\nu_{120}\) et conserve \(\nu_{270}\). Cette parité
définit l'architecture historique d'un lecteur indépendant des grandeurs
métrologiques utilisées comme cible. Son application
littérale au registre vivant \(81\times108\), en prenant comme
\(\varepsilon(t)\) l'unique bit global de signe conservé par cette
projection, produit
\[
 s_C^{(m)}(\nu)=0,
 \qquad \tau_m^\nu=\nu_{120}=\nu_{270}=0
\]
aux quatre-vingt-une positions. Les événements antipodaux ont été
identifiés par la projection du registre. Le résultat local est conservé
comme lecteur récupéré et, en même temps, comme test de perte d'information :
l'extracteur neutre de
\cref{sec:extractor-neutro-z0-torsion-rev2} agit avant cette
identification et lit le bit bidirectionnel local, le feuillet, la torsion et la
couronne antipodale. Le zéro obtenu après projection ne dégrade pas cette
construction ; il certifie exactement quelle information le lecteur global efface.
Les référentiels
\(F_\ell,F_\nu\), le lecteur~\eqref{eq:invariantes-neutrino-ledger-autonoma}
et la forme agrégée
\begin{equation}
 G_{ab}=\frac1{\sqrt{|L_a||N_b|}}
 \sum_{c\in L_a}\sum_{d\in N_b}K_{\rm registro}(c,d),
 \qquad U_{\rm int}=\operatorname{polar}(G),
 \label{eq:kernel-ledger-pmns-autonoma}
\end{equation}
fixent ainsi les deux extrémités déjà construites du pont PMNS.

L'égalité cardinale entre les cinquante-six sommets multiplicatifs des
germes et les cinquante-six familles CT108 a été soumise à un premier
contrôle indépendant des grandeurs métrologiques utilisées comme cible.
L'application de Hadamard naturelle sur un seul germe n'atteint
que vingt-neuf familles ; en parcourant les symétries diédriques, les signes et les
translations naturelles, la couverture maximale est de trente-cinq. Par conséquent,
l'entrelaceur agit sur le couple de germes enrichis et conserve feuillet,
retenue et mémoire. La matrice d'incidence comparant les émissions partagées
dans ce domaine double devra être de rang cinquante-six, être naturelle sous
orientation et ne pas se factoriser par des compteurs globaux. Ce critère fixe le
domaine exact de la construction suivante.

\endgroup
\begingroup
\let\chapter\section
\let\section\subsection
\let\subsection\subsubsection
\let\subsubsection\paragraph
\let\clearpage\relax
\chapter{Algèbre qutrit de couleur et connexion finie de jauge}
\label{chap:algebra-color-qutrit}
\index{color!qutrit}\index{opérateurs de Weyl}
\index{groupe spécial unitaire}\index{action de Wilson}

La fonctionnelle résiduelle \(\kappa_3\) du
\cref{chap:ocupacion-estadistica} fournit trois classes affines ; l'atlas
ultérieur réalise leur bilan \(12+12+12\) dans la
\cref{prop:color-atlas}. Un triplet d'étiquettes n'est pas encore une
représentation. L'étape algébrique requiert de déclarer la carte qutrit, son
orientation et son caractère de phase. Sous ces hypothèses, la représentation et
son algèbre infinitésimale s'obtiennent exactement, sans utiliser de masses ni de données
chromodynamiques.

\section{Couple de Weyl et algèbre de trace \(0\)}
\label{sec:weyl-color-qutrit}

Soit
\[
 V_{\rm col}=\mathbb C^3,
 \qquad
 \omega=e^{2\pi i/3},
\]
avec la base ordonnée \((e_0,e_1,e_2)\). Nous définissons
\[
 Xe_j=e_{j+1\bmod3},
 \qquad
 Ze_j=\omega^j e_j.
\]
On a
\[
 X^3=Z^3=I,
 \qquad
 ZX=\omega XZ.
\]

\begin{theorem}[Base de Weyl et algèbre complexe de couleur]
\label{thm:weyl-sl3-color}
Pour
\[
 W_{ab}:=X^aZ^b,
 \qquad
 (a,b)\in(\mathbb Z/3\mathbb Z)^2,
\]
les neuf opérateurs \(W_{ab}\) forment une base orthogonale de
\(\operatorname{End}_{\mathbb C}(V_{\rm col})\) pour le produit de
Hilbert--Schmidt. Les huit opérateurs différents de l'identité ont une trace nulle et
\[
 \operatorname{span}_{\mathbb C}
 \{W_{ab}:(a,b)\ne(0,0)\}
 =\mathfrak{sl}_3(\mathbb C).
\]
\end{theorem}

\begin{proof}
La relation de Weyl permet de réduire tout produit à la forme
\(X^aZ^b\). Un calcul direct donne
\[
 \operatorname{tr}(W_{ab}^{\dagger}W_{cd})
 =3\,\delta_{ac}\delta_{bd}.
\]
Par conséquent, les neuf opérateurs sont linéairement indépendants ; comme
\(\dim_{\mathbb C}\operatorname{End}_{\mathbb C}(V_{\rm col})=9\), ils forment une
base. La trace de \(X^aZ^b\) s'annule sauf pour
\((a,b)=(0,0)\). Les huit autres forment ainsi une base du sous-espace de
trace nulle, dont la dimension complexe est huit.
\end{proof}

\begin{corollary}[Forme réelle compacte]
\label{cor:su3-color-qutrit}
L'involution
\[
 \tau(A):=-A^\dagger
\]
préserve \(\mathfrak{sl}_3(\mathbb C)\), et son ensemble de points fixes est
\[
 \mathfrak{sl}_3(\mathbb C)^\tau
 =\{A:A^\dagger=-A,\ \operatorname{tr}A=0\}
 =\mathfrak{su}(3).
\]
Cette forme réelle est de dimension huit.
\end{corollary}

\begin{proof}
On a
\[
 W_{ab}^{\dagger}=\omega^{ab}W_{-a,-b}.
\]
Les huit indices non nuls forment quatre couples adjoints. Pour
\[
 \mathcal R=\{(1,0),(0,1),(1,1),(1,2)\},
\]
les opérateurs
\[
 A_u=W_u-W_u^\dagger,
 \qquad
 B_u=i(W_u+W_u^\dagger),
 \qquad u\in\mathcal R,
\]
sont antihermitiens, de trace nulle et linéairement indépendants sur
\(\mathbb R\). Ils constituent donc une base réelle de
\(\mathfrak{su}(3)\).
\end{proof}

L'application résiduelle
\[
 \kappa_{\rm col}(r,c)
 =\mathbb C e_{\,r-c\bmod3}
\]
associe les trois classes de l'atlas aux trois droites propres de \(Z\).
L'opérateur \(X\) les permute cycliquement et \(Z\) enregistre leur phase. Cette
attribution est exacte une fois fixés
\[
 V_{\rm col}=\mathbb C[\mathbb F_3],
 \quad
 \text{la base delta},
 \quad
 j\longmapsto j+1,
 \quad
 j\longmapsto\omega^j.
\]
Les changements affines admissibles de la carte produisent des représentations
unitairement conjuguées. La construction est donc naturelle à
conjugaison unitaire près ; elle ne sélectionne pas à elle seule les noms physiques des
couleurs ni une réalisation QCD.

\section{Connexion et action sur un complexe cellulaire fini}
\label{sec:gauge-finito-color}

\begin{construction}[Théorie de jauge finie relative au complexe]
\label{cons:gauge-finito-color}
Soit \(\mathcal K=(V,E,F)\) une cellularisation finie orientée compatible avec
le graphe des chemins. À chaque arête orientée \(e\), on attribue
\[
 U_e\in SU(3),
 \qquad
 U_{\bar e}=U_e^{-1}.
\]
Une transformation \(g:V\to SU(3)\) agit par
\[
 U_e\longmapsto U_e^g
 =g_{t(e)}U_eg_{s(e)}^{-1}.
\]
Pour chaque face \(f\), le produit ordonné
\[
 U_f=\prod_{e\subset\partial f}U_e^{\epsilon(f,e)}
\]
définit son holonomie. Un sommet de base \(v_f\) étant choisi, on a
\[
 U_f^g=g_{v_f}U_fg_{v_f}^{-1}.
\]
\end{construction}

\begin{proposition}[Invariance de jauge et transport covariant]
\label{prop:wilson-color-finito}
Pour \(\beta\ge0\), l'action
\[
 S_{\rm col}^{(\mathcal K,\beta)}[U]
 =\frac{\beta}{3}\sum_{f\in F}
 \left(3-\operatorname{Re}\operatorname{tr}U_f\right)
\]
est invariante de jauge et non négative. Si
\(\psi_v\in V_{\rm col}\) se transforme comme
\(\psi_v\mapsto g_v\psi_v\), alors
\[
 \nabla_e\psi:=U_e\psi_{s(e)}-\psi_{t(e)}
\]
satisfait
\[
 \nabla_e^g\psi^g=g_{t(e)}\nabla_e\psi.
\]
\end{proposition}

\begin{proof}
L'holonomie change par conjugaison et la trace est cyclique ; d'où
l'invariance. Comme chaque \(U_f\) est unitaire de dimension trois,
\(\operatorname{Re}\operatorname{tr}U_f\le3\), ce qui prouve la
non-négativité. La covariance de \(\nabla_e\) résulte de la substitution des deux lois
de transformation.
\end{proof}

\section{Positivité des noyaux de classe et décomposition harmonique}

La connexion finie précédente n'implique pas à elle seule la positivité par
réflexion ni une limite de Yang--Mills. Avant de formuler de tels transferts,
il convient d'isoler le contrôle classique qui est effectivement clos. Soit \(G\) un groupe
compact et \(\widehat G\) son dual unitaire. Pour chaque
\(\rho\in\widehat G\), soient \(d_\rho\) sa dimension et
\(c_\rho(a)\ge0\) des coefficients de décroissance suffisante, par rapport à la
croissance de \(d_\rho\), pour que la série suivante converge
absolument.

\begin{definition}[Noyau de classe]
\label{pf84:YM-07}
On définit
\[
 K_a(g,h):=
 \sum_{\rho\in\widehat G}
 c_\rho(a)\,
 \Tr\!\bigl(\rho(g)\rho(h)^*\bigr),
 \qquad g,h\in G.
\]
Le choix
\[
 c_\rho(t)=d_\rho e^{-tC_2(\rho)},
 \qquad t>0,
\]
fournit l'exemple du noyau de chaleur. La non-négativité des
coefficients ne remplace pas l'hypothèse de convergence absolue.
\end{definition}

\begin{theorem}[Positivité de Peter--Weyl]
\label{pf84:YM-08} Pour \(g_1,\dots,g_n\in G\) et \(z_1,\dots,z_n\in\CC\),
\[
 \sum_{i,j=1}^n
 \overline z_i z_jK_a(g_i,g_j)\ge0.
\]
\end{theorem}

\begin{proof}
L'unitarité des représentations et la convergence absolue permettent
de réordonner la somme :
\[
 \sum_{i,j}\overline z_i z_jK_a(g_i,g_j)
 =
 \sum_{\rho\in\widehat G}
 c_\rho(a)
 \left\lVert
 \sum_j\overline z_j\rho(g_j)
 \right\rVert_{\mathrm{HS}}^2\ge0.
\]
Un coefficient spectral négatif peut détruire cette positivité de Gram ;
le signe fait donc partie des hypothèses.
\end{proof}

\begin{lemma}[Contrôle abélien de Poisson]
\label{pf84:YM-09}
Pour \(G=U(1)\), \(\theta\in\RR\) et \(|a|<1\),
\[
 \sum_{n\in\ZZ}a^{|n|}e^{in\theta}
 =\frac{1-a^2}{1-2a\cos\theta+a^2}.
\]
Si l'on exige en outre des coefficients de Fourier non négatifs, on restreint
\(a\) à \(0\le a<1\).
\end{lemma}

\begin{proof}
L'égalité résulte de la sommation séparée des deux séries géométriques
d'indices positifs et négatifs et de l'ajout du mode \(n=0\). Pour \(a<0\),
l'identité scalaire reste valable, mais ne constitue plus un exemple de
coefficients de Peter--Weyl non négatifs.
\end{proof}

\begin{remark}
Les \cref{pf84:YM-07,pf84:YM-08,pf84:YM-09} sont des définitions et des résultats
classiques d'analyse harmonique. Ils fournissent un contrôle positif pour
les noyaux de classe ; ils ne démontrent ni la positivité d'Osterwalder--Schrader, ni l'écart de
masse, ni le confinement, ni la limite continue, et ne promeuvent pas la théorie de jauge
finie de ce chapitre au rang de chromodynamique quantique.
\end{remark}

\begin{exactbox}[title={Clôture algébrique de la couleur}]
Les \cref{thm:weyl-sl3-color,cor:su3-color-qutrit} construisent exactement la
chaîne
\[
\text{reste ternaire}\longrightarrow
\text{qutrit}\longrightarrow
\mathfrak{sl}_3(\mathbb C)\longrightarrow\mathfrak{su}(3).
\]
La \cref{cons:gauge-finito-color,prop:wilson-color-finito} prolonge cette
chaîne en une connexion et une action de jauge sur la cellularisation finie
\(\mathcal K\). Tel est l'objet démontré dans le chapitre. La
comparaison avec la chromodynamique physique s'effectue ensuite et ne redéfinit ni le
qutrit, ni le bilan \(12+12+12\), ni la représentation construite.
\end{exactbox}

\endgroup
\begingroup
\let\chapter\section
\let\section\subsection
\let\subsection\subsubsection
\let\subsubsection\paragraph
\let\clearpage\relax
% Vista editorial de corpus/source/masas/02_ley_general.tex, líneas 363--491: cuerpo exacto salvo el retítulo académico del epígrafe sobre el ciclo de 108 estados. Las líneas 1--263 ya residen exactamente en 72_rev2_nucleo_ley_i.tex y las líneas 264--362 en 72_rev2_algoritmo_universal.tex.
\chapter{Extension systématique de la théorie des masses}
\Needspace{7\baselineskip}
\section{Fondement unifié de la masse}
\Needspace{7\baselineskip}
\subsection{L'objet antérieur à la grandeur}
La théorie commence avant la valeur numérique. Une masse physiquement individualisée n'est pas un nombre réel accompagné du nom d'une particule, mais le terme final d'une biographie mathématique. APP fournit l'espace fini d'incidences et les deux opérations corrélatives ; le TRIT conserve l'orientation temporelle de chaque transition ; le TPK transporte phase, feuillet, retenue, frontière, parcours, mémoire et survie. La Mécanique Dimensionnelle convertit ensuite cette biographie en une section de la droite de masse. L'ordre complet est

\[
\mathrm{APP}\longrightarrow\mathrm{TRIT}\longrightarrow\mathrm{TPK}
\longrightarrow\Gamma^{\rm surv}
\longrightarrow\gamma_{\rm obs}
\longrightarrow\mathscr L_\gamma
\longrightarrow\sigma(\gamma)
\longrightarrow\chi_{\rm mass}
\longrightarrow\mathcal T_{90,120}
\longrightarrow\mathcal L_M.
\]
Chaque flèche retient une information utilisée par la suivante. La cellule ne remplace pas le parcours ; le parcours ne remplace pas l'enregistrement ; une signature ne remplace pas son enregistrement généalogique ; un caractère ne remplace pas la section dimensionnelle sur laquelle il agit. Cette séparation explique pourquoi deux particules peuvent partager charge ou spin tout en possédant des masses différentes : leur biographie enrichie n'est pas la même. Elle explique aussi pourquoi particule et antiparticule ont la même masse : l'involution change l'orientation impaire et conserve la mémoire paire qui descend au quotient de masse.

APP comprend neuf fois neuf positions publiées. Chaque position admet quatre orientations initiales, si bien que le domaine orienté contient

\[
81\cdot4=324
\]
parcours candidats. La classification interne produit 56 familles de parcours et 13 multisections. Ces cardinalités sont simultanées : 81 compte les cellules, 324 les orientations, 56 les quotients par histoire et 13 les types opératoires grossiers. Aucun tableau de particules ne remplace ces domaines.

\Needspace{7\baselineskip}
\subsection{Unité discrète du continu qui soutient la masse}
La droite de masse ne s'adosse pas à cinq théories indépendantes du continu. Elle procède de l'unique état APP--TRIT--TPK dont émergent corrélativement les aspects spectral, solénoïdal, de jauge, cohérent et déterminantal, désignés dans la notation de construction par \(sp,sol,gau,coh,det\). Ces symboles sont des discriminants de lecture, et non des sous-domaines préalables. Si \(\mathcal C_{\rm disc}\) désigne la structure commune, la flèche génératrice prend la forme

\[
\mathfrak E_{\rm cont}:\mathcal C_{\rm disc}
\longrightarrow
\bigl(\mathcal C_\infty;
\mathcal O_{sp},\mathcal O_{sol},\mathcal O_{gau},
\mathcal O_{coh},\mathcal O_{det}\bigr),
\]
où les cinq observations appartiennent à un même objet limite, partagent fibres, orientation, retenue, mémoire, frontière et lecteurs, et sont transportées par les mêmes carrés de raffinement. La masse utilise sa lecture déterminantale pour l'échelle d'action, sa lecture spectrale pour le caractère et les tours, sa lecture de jauge pour les secteurs nuls et chargés, et ses lectures cohérente et solénoïdale pour la propagation et la composition. Aucune n'est introduite en supposant d'emblée que l'état appartient à une classe physique classique.

Cette unité est une condition de typage : séparer l'un de ces aspects en conservant les autres produirait un objet distinct, incapable de soutenir simultanément parcours, action, charge, amplitude et masse. La présente théorie ne redémontre pas le continu complet ; elle en incorpore la construction unitaire comme fondement et vérifie qu'aucune des insertions relatives aux masses ne le dissocie.

\Needspace{7\baselineskip}
\subsection{Échelle absolue et quotients}
Le lecteur déterminantal local a pour forme normale de Smith

\[
\operatorname{SNF}(H_4)=\operatorname{diag}(1,2,2,4),
\]
d'où

\[
|\operatorname{coker}H_4|=16,
\qquad
\Sigma_H=\varphi_{\rm HMT}\alpha_{\rm HMT}^{16},
\qquad
\operatorname{ord}_{10}(\Sigma_H)=-34.
\]
La décennie est obtenue avant la carte SI. Les sections pré et de retour de la droite d'action sont

\[
\hbar_{\rm pre}=H_5\,10^{-34}u_S,
\qquad
\hbar_{\rm ret}=\eta_{\rm ret}^{(\deg)}\,10^{-34}u_S,
\]
et leur quotient bidirectionnel est

\[
R_{\rm act}=\frac{H_5}{\eta_{\rm ret}^{(\deg)}}.
\]
Dans la branche d'horloge, nous choisissons explicitement comme section d'action la section de retour, \(\hbar_{\rm int}:=\hbar_{\rm ret}\). Une horloge déclarée \(t_0\) fournit

\[
\omega_0=\frac{2\pi}{108t_0},
\qquad
\varepsilon_{\rm clk}=\hbar_{\rm ret}\omega_0,
\qquad
m_{\rm clk}=\varepsilon_{\rm clk}c_{\rm int}^{-2}.
\]
Adopter \(u_M^{\rm clk}:=m_{\rm clk}\) comme origine de la carte constitue une normalisation explicite. Cela n'identifie pas automatiquement cette section à une base énergétique ou massique préalablement déclarée. Le caractère produit d'abord l'opérateur de base

\[
\widehat M_X^{(0)}
=\mathfrak e_0m_{\rm clk}\,
\chi_{\rm full}(\sigma_X-\sigma_e,\eta_X-\eta_e),
\]
et le lecteur bidirectionnel agit ensuite. Pour chaque présentation \(r\in\{D,\Omega\}\), soit \(\widehat B_X^r=(\widehat B_X^r)^*\) sur la fibre finie ; dans une réalisation infinie, un domaine commun dense est en outre exigé pour le calcul fonctionnel. Alors

\[
\widehat M_X^{\beta,r}
=R_{\rm act}^{\widehat B_X^r/2}
\widehat M_X^{(0)}
R_{\rm act}^{\widehat B_X^r/2}.
\]
Dans une fibre scalaire, \(\widehat B_X^r=\kappa_XI\), de sorte que

\[
m_X^\beta=m_X^{(0)}R_{\rm act}^{\kappa_X}.
\]
Si les coordonnées de tour sont déjà exprimées relativement à l'électron, on omet \(-\eta_e\). Dans un quotient, la section dimensionnelle commune se simplifie, mais non la différence de lecture bidirectionnelle :

\[
\frac{m_Y^\beta}{m_X^\beta}
=\chi_{\rm full}(\sigma_Y-\sigma_X,\eta_Y-\eta_X)
R_{\rm act}^{\kappa_Y-\kappa_X}.
\]
Ainsi, \(10^{-34}\) fixe l'échelle absolue d'action, d'énergie et de masse ; ce n'est pas un correcteur relatif caché. L'information fine d'une espèce ne peut intervenir que par son caractère, son horloge relative ou le quotient orienté des sections pré et de retour.

\Needspace{7\baselineskip}
\subsection{Quantum interne de la droite de masse induit par le cycle de 108 états}
L'autodualité Planck--CT108 fournit, dans la normalisation interne

\[
\hbar_{\rm int}=c_{\rm int}=t_0=1,
\]
la section

\[
m_{108}^{\rm int}
=0.058177641733144319230789692282953757\ldots .
\]
Ce nombre est un quantum normalisé de la droite interne de masse. Il n'est pas la masse d'une particule, ne porte pas à lui seul une carte MeV ou GeV et ne peut être inséré comme ligne supplémentaire d'espèces. Il sert à vérifier la compatibilité entre la carte d'action, l'horloge CT108 et l'autodualité de la droite de masse. La transition vers une masse nominale exige encore le parcours, la signature, l'opérateur et la section dimensionnelle de l'espèce.

\Needspace{7\baselineskip}
\subsection{Existence du résultat HMT–MD pour les extensions survivantes et critère de réduction scalaire selon le rang de la fibre}
Une extension survivante, un lecteur cellulaire, une signature entière, une famille absolument convergente de coefficients de tour et une section dimensionnelle de masse étant fixés, il existe un résultat HMT--MD correctement typé. Si la fibre possède un unique point fixe compatible avec l'involution, le résultat se réduit à un scalaire canonique. Si la fibre est de rang supérieur, le résultat primaire est l'opérateur spectral sur la fibre ; une masse scalaire physique requiert la fonctionnelle de couplage qui individualise la représentation observée.

Ce théorème ne déprécie pas les évaluations scalaires retrouvées. Il distingue trois affirmations : l'opérateur interne est construit ; l'évaluation d'une signature déclarée est exacte ; la détermination prospective de cette signature constitue un théorème supplémentaire lorsqu'elle ne se déduit pas du lecteur lui-même. La comparaison expérimentale n'est effectuée qu'après avoir fixé laquelle de ces trois affirmations est présentée.

\Needspace{7\baselineskip}

\endgroup

\section{Algorithme universel de réalisation matérielle et égalisateur spectral des lecteurs \(K_D\) et \(K_\Omega\)}
La construction s’achève par l’algorithme qui reçoit uniquement des données généalogiques typées, forme la fibre de réalisation et exige l’égalisateur des lecteurs avant de produire une sortie dimensionnelle.

\begingroup
\let\chapter\section
\let\section\subsection
\let\subsection\subsubsection
\let\subsubsection\paragraph
\let\clearpage\relax
\section{Algorithme universel de la Loi Générale des Masses}
\Needspace{7\baselineskip}
\subsection{Domaine du lecteur cellulaire : extension survivante, parcours, enregistrement et signature}
L'entrée contient uniquement :

\[
 (\mathrm{APP},\mathrm{TRIT},\mathrm{TPK},
 \mathcal R_{324},\mathcal L_{108},
 \text{représentation},\text{couplage},
 \text{carte dimensionnelle commune}).
\]
Elle ne contient ni masse observée, ni résidu, ni coefficient choisi par approximation, ni tableau externe.

\Needspace{7\baselineskip}
\subsection{Construction fonctorielle de la réalisation matérielle : fibre de parcours, entrelaceur, projecteur de coïncidence, opérateur spectral et scalarisation}
Pour chaque état physique \(X\) :

\begin{enumerate}
\item construire son descripteur de charge, génération, saveur, couleur, spin, statistique,
orientation et composition ;

\item former la fibre des parcours compatibles ;
\item calculer l'espace des entrelaceurs
\end{enumerate}
   \[
   \mathscr I_X=
   \operatorname{Hom}_{G_X}(V_X,\ell^2(F_X));
   \]
\begin{enumerate}
\item si \(\dim\mathscr I_X=0\), écarter cette identification ;
\item construire le projecteur par moyenne de Reynolds et facteur polaire ;
\item reconstruire les enregistrements généalogiques de tous les parcours du support ;
\item calculer la signature, \(K_D\), \(K_\Omega\) et
\(\mathcal F_{\rm tower}^{2w}\) ;

\item appliquer le pointeur de lecture du couplage, l'égalisateur ou la mesure
spectrale conjointe ;

\item construire \(\widehat M_X\) ;
\item réduire à un scalaire seulement si la variance est nulle ou s'il existe un pôle isolé ;
\item pour les composés, appliquer d'abord \(\mathfrak C_{12}\) ;
\item pour les largeurs, calculer le semi-groupe de survie ;
\item multiplier par la section dimensionnelle commune ;
\item figer le résultat ;
\item ouvrir la comparaison externe.
\end{enumerate}
\Needspace{7\baselineskip}
\subsection{Existence de la réalisation par fibres}
La multiplicité d'une représentation \(\pi_X\) au sein d'une fibre se calcule par

\[
 \dim\mathscr I_X=
 \frac1{|G_X|}
 \sum_{g\in G_X}
 \overline{\chi_{\pi_X}(g)}
 |\operatorname{Fix}_{F_X}(g)|.
\]
Si elle est positive, il existe un entrelaceur. Si elle est supérieure à un, le résultat est un faisceau de réalisations et le couplage doit sélectionner un état ou conserver la mesure mixte. Cette formule remplace une recherche nominale par une condition algébrique finie sur les 324 parcours.

\Needspace{7\baselineskip}
\subsection{Égalisateur spectral des lecteurs \(K_D\) et \(K_\Omega\) au moyen de \(Q_{\mathrm{eq}}=\mathbf1_{\{0\}}(\widehat L_D-\widehat L_\Omega)\)}
Soient

\[
 L_r(\Gamma)=
 \kappa_r(\Gamma)\log R_{\rm act}
 +\sum_{h\in\mathfrak S}\eta_h^r(\Gamma)s_h,
 \qquad r\in\{D,\Omega\}.
\]
Le projecteur de coïncidence est

\[
 Q_{\rm eq}=
 \mathbf1_{\{0\}}(\widehat L_D-\widehat L_\Omega).
\]
Une réalisation admet une lecture commune exactement lorsque

\[
 P_X\le Q_{\rm eq}.
\]
S'il existe une involution du couplage qui échange les deux lecteurs, on utilise \(L_{\rm sym}\). Sinon, on présente leur mesure spectrale conjointe.

\Needspace{7\baselineskip}
\subsection{Codomaine du lecteur cellulaire : caractère adimensionnel, tour et réalisation de masse}
Le résultat peut être :

\[
 \begin{cases}
 m_X\in\mathcal L_M,&\text{masse scalaire},\\
 \operatorname{Spec}(\widehat M_X),&\text{multiplet},\\
 z_X=M_X-i\Gamma_X/2,&\text{résonance},\\
 0\in\mathcal L_M,&\text{section nulle},\\
 (\mathcal C,F,R,\Delta E),&\text{secteur topologique},\\
 (L,\operatorname{Obs}_L),&\text{section virtuelle}.
 \end{cases}
\]
Cette diversité de codomaines fait partie de la loi ; elle n'est pas une exception.

\Needspace{7\baselineskip}

\endgroup

\section{Bifurcation nulle et positivité}
Le secteur photonique et gluonique est séparé avant l’application du caractère positif. L’identité \(R_{\rm act}^{0}=1\) ne démontre pas une masse nulle ; la nullité provient de la carte de jauge. Dans les familles matérielles, la sortie primaire est un opérateur positif de fibre. Seule une compression physique typée produit un scalaire.

\section{Accès au domaine exhaustif}
La définition et les théorèmes du domaine précèdent le déploiement de ses
\hyperref[chap:atlas-genealogico-exhaustivo-324-rutas]{324 fichas completas}. Les fiches sont placées à la fin de l’appareil scientifique afin de conserver la continuité de la dérivation sans dissimuler aucun parcours.
